% ============================================================
%  第二卷 第十章 引力场中的粒子（§81–§90 + 习题）
%  底本：高教社中文版第8版（鲁欣/任朗 译）；OCR 错误已修正，
%  脚注文本待从纸版补录（正文圈码 ①②③ 已按原书保留）
%  已修正的 OCR/原书错误：
%   - §81 "称为引力势 φ"：OCR 语序错乱（"其中，是关于坐标…称为引力势 φ①"），
%     按"其中 φ 是关于坐标和时间的某一函数……称为引力势①"复原
%   - §82 (82.2)：OCR ".stackrel.per.i≠k" 乱码 → "当 i≠k"
%   - §82 "时空的几何性质"一句 OCR 行错位（"由gi各量"），已复原
%   - §83 (83.5)(83.6)(83.7)：变换公式分母哑标错乱（∂x^n→∂x'^l 等），按张量变换法则复原
%   - §83 dΩ'→(1/J)dΩ：OCR 丢失箭头，按原书补 →
%   - §84 "引入三维矢量 g"：OCR 误作 ε/𝐄，按原书 p263 为三维矢量 g
%   - §84 脚注①（二次型 (84.6) 正定性条件、g_ik 行列式不等式组）整段为脚注文本，随脚注代理补录
%   - §85–§86：Γ（克里斯托夫符号）OCR 大量误作 T/I/cal T，统一为 \Gamma；
%     "DA^i" 的 unicode 组合字符 Ḍ → \mathrm{D}；δ(A_iB^i)=0 的 δ OCR 误作 S/8
%   - §86 (86.1) 前：g_{ik;l} 展开式中间项 OCR "T_{ll}^m" → Γ^m_{il}
%   - §87 拉格朗日量 (87.10) 前常数 "-mc^2" 后附圈码②：OCR 污染为 "-mc^{2②}"，已分离
%   - §87 习题：δ(g_ik dx^i dx^k) 的 δ 误作 8；dδx^k/ds 的分母 OCR 误作 ∂s
%   - §88 习题2 标题"费马原理"：OCR 误作"费 π/2 原理"
%   - §88 习题1 (1)：Γ 分量表达式三条按原书 p280 逐字复原（OCR 乱码）
%   - §89 习题"利用（84.6），（84.7）"：原书/OCR 均印作（86.4），（86.7），
%     所引实为三维度规 (84.6)(84.7)，按物理修正（原书互参排印错误）
%   - §90 正文中的协变形式 F_{ik;l}+F_{li;k}+F_{kl;i}=0 为脚注①文本，剔出正文
%  遗留问题：
%   - §87 脚注①含公式 (87.3a)（四维加速度形式的运动方程），随脚注文本一并留待脚注代理
%   - §88 习题1 脚注①（三维单位反对称张量 η_αβγ 及 rot/div 定义）待脚注代理
%  脚注已转录（25 处圈码全部替换为 \footnote，文本逐页对照原书第 251–284 页核校）
% ============================================================
\chapter{引力场中的粒子}

\section{非相对论力学中的引力场}

引力场（或者重力场）具有以下的基本特性：所有的物体，不管质量的大小，只要有相同的初始条件，它们在场中就将以相同的方式运动.

例如，对于在地球引力场中的所有物体，自由下落的规律都是一样的，即不管它们的质量如何，都获得同一的加速度.

引力场的这个特性使我们有可能确定以下两种运动有类似之处：一种是一个物体在引力场中的运动；另一种是一个物体不在任何外场中的运动，但是在一个非惯性参考系中考察.实际上，在一个惯性参考系中，所有物体的自由运动都是匀速直线运动；假如在开始时，它们的速度是一样的，那么无论在任何时候都将保持一样.因此，很显然，假如我们考虑在一给定的非惯性系中的自由运动，那么，相对于这个参考系，所有的物体都将以同样的方式运动.

因此，在一个非惯性系中的运动特性与在一个有引力场存在的惯性系中的运动特性一样；换句话说，非惯性系与某一引力场等效.这种情况称为等效原理.

我们来研究，例如，匀加速参考系中的运动.假如一个任意质量的物体在这样的参考系中作自由运动，很显然，该物体对于这个参考系就有一个恒定的加速度，这个加速度与参考系本身的加速度大小相等，方向相反.这样的描述也适用于物体在均匀恒定引力场，例如地球的引力场中的运动（在不大的区域内，地球的引力场可以当做是均匀的）.因此，匀加速参考系与一个不变的、恒定的外场等效.同理，参考系的非匀加速直线运动，显然与一个均匀但变化着的引力场等效.

然而必须着重指出，与非惯性参考系等效的场其实并不完全与“实际的”引力场一样（后者在惯性系中也存在），例如，它们在无穷远处的性质就有本质的区别.在与产生场的物体相距为无穷远处，“实际的”引力场总是趋近于零；与此相反，与非惯性参考系等效的场在无穷远处却无限制地增大，或者，无论如何，总保持为有限值.例如，在旋转参考系中出现的离心力，当我们从旋转轴离开时，无限制地增大；一个与作匀加速直线运动的参考系等效的场在空间各处都是一样的，而且在无穷远处也一样.

只要我们从非惯性系过渡到惯性系，与非惯性系等效的场就消失了.与此相反，无论选择哪一种参考系，“实际的”引力场（在惯性参考系内也存在）总是无法消除的.这可以从上面所讲的关于“实际的”引力场和与非惯性系等效的场在无穷远处情况的差别直接看出，既然后者在无穷远处不趋近于零，那么，无论怎样选择参考系，“实际的”场也不可能消除，因为它在无穷远处变为零.

选择适当参考系的方法只能消除空间中某一给定区域内的引力场，而且这个区域必须如此之小，使其中的场可以看做是均匀的.要做到这点，可选择一个作加速运动的参考系，这个加速度就等于粒子放在我们正在考虑的场的区域内所得到的加速度.

粒子在引力场中的运动，在非相对论力学中为拉格朗日量所决定，它在惯性参考系中的表达式是
\begin{equation}
  L=\frac{mv^2}{2}-m\varphi,
\end{equation}
其中 $\varphi$ 是关于坐标和时间的某一函数，它可描述场的特性，称为引力势\footnote{以后我们不需用电磁势 $\varphi$，因此用同一个符号 $\varphi$ 来代表引力势不会引起误解.}.与此相应，粒子的运动方程是
\begin{equation}
  \dot{\boldsymbol{v}}=-\operatorname{grad}\varphi.
\end{equation}
这个方程不包含质量，也不包含任何其他描述粒子特征的常数，它就是引力场基本特性的数学表示.

\section{相对论力学中的引力场}

在上节中所指出的引力场的基本特性，即所有物体在场中作同样的运动，在相对论力学中也还是有效的.所以，引力场和非惯性参考系的类似性依然存在.因此，在研究相对论力学中引力场的特性时，我们自然也从这个类似性出发.

在惯性参考系中，用笛卡儿坐标，间隔 ds 由关系式
\begin{equation*}
  \mathrm{d}s^2=c^2\mathrm{d}t^2-\mathrm{d}x^2-\mathrm{d}y^2-\mathrm{d}z^2
\end{equation*}
给定.在变换到任何其他惯性参考系时（就是说在作洛伦兹变换时），我们知道，间隔保持同样的形式.然而，假如我们变换到非惯性系，$\mathrm{d}s^2$ 将不再是四个坐标的微分的平方和了.

例如，当我们变换到匀速旋转的坐标系时，
\begin{equation*}
  x=x'\cos\varOmega t-y'\sin\varOmega t,\qquad
  y=x'\sin\varOmega t+y'\cos\varOmega t,\qquad
  z=z'
\end{equation*}
（Ω 是旋转的角速度，其方向沿着 z 轴），那么，间隔就有下面的形式：
\begin{equation*}
\begin{aligned}
  \mathrm{d}s^2={}&\left[c^2-\varOmega^2(x'^2+y'^2)\right]\mathrm{d}t^2
  -\mathrm{d}x'^2-\mathrm{d}y'^2\\
  &-\mathrm{d}z'^2+2\varOmega y'\mathrm{d}x'\mathrm{d}t
  -2\varOmega x'\mathrm{d}y'\mathrm{d}t.
\end{aligned}
\end{equation*}
不管时间坐标变换的规律怎样，这个式子不可能仍以四个坐标的微分的平方和来表示.

因此，在非惯性参考系中，间隔的平方是坐标微分的一般形式的二次型，也就是说，它有以下的形式：
\begin{equation}
  \mathrm{d}s^2=g_{ik}\mathrm{d}x^i\mathrm{d}x^k,
\end{equation}
其中 $g_{ik}$ 是空间坐标 $x^1$，$x^2$，$x^3$ 和时间坐标 $x^0$ 的某些函数.因此，假如我们使用非惯性系，四维坐标 $x^0$，$x^1$，$x^2$，$x^3$ 将是曲线坐标.$g_{ik}$ 可用来表示时空度规，它决定每个曲线坐标系的所有几何特性.

显然，总可以认为 $g_{ik}$ 这些量对指标 i 和 k 来说是对称的（即 $g_{ik}=g_{ki}$），因为它们是由对称式（82.1）所决定的，此处的 $g_{ik}$，$g_{ki}$ 是作为同一个积 $\mathrm{d}x^i\mathrm{d}x^k$ 的因子出现的.在一般情形下，一共有十个不同的量 $g_{ik}$，四个有相同的指标和 $4\times3/2=6$ 个有不同的指标.在一个惯性参考系中，当我们用笛卡儿空间坐标 $x^{1,2,3}=x,y,z$ 和时间坐标 $x^0=ct$ 时，$g_{ik}$ 各量是
\begin{equation}
  g_{00}=1,\quad g_{11}=g_{22}=g_{33}=-1,\quad
  g_{ik}=0\quad\text{当}\ i\neq k.
\end{equation}
具有这些 $g_{ik}$ 值的四维坐标系称为伽利略坐标.

在上一节中已经证明了非惯性参考系与某些力场等效.现在我们看出，在相对论力学中，这些场由 $g_{ik}$ 各量所决定.

这同样也可以应用到“实际的”引力场.任何引力场都只是时空度规的一个改变，而这个改变是由 $g_{ik}$ 各量所决定的.这个重要事实说明，时空的几何性质（它的度规）是由物理现象所决定的，而不是空间和时间的固定不变的性质.

建立在相对论基础上的引力场理论称为广义相对论.它是由爱因斯坦提出来的（最后在1915年建立），在现有的物理理论中，它或许是最美丽的.突出之处是爱因斯坦用纯演绎的方法就建立了这个理论，只是在以后才被天文观测所证实.

正如非相对论力学中一样，在“实际的”引力场和与非惯性参考系等效的场之间有一个根本的差别.当变换到非惯性参考系时，二次型具有(82.1)的形式，即 $g_{ik}$ 各量，是通过简单的坐标变换从伽利略坐标的值（82.2）得到的，所以，可利用逆坐标变换在整个空间中回到伽利略坐标的值.这种形式的 $g_{ik}$ 非常特殊，这是因为在一般情况下仅仅只用四个坐标的变换就要使 $g_{ik}$ 的十个量化为预定的形式是不可能的.

“实际的”引力场用任何坐标变换都不能消除；换句话说，存在引力场时，时空是这样的，不可能用任何坐标变换使决定其度规的量 $g_{ik}$ 在整个时空内都化为伽利略坐标的值.这样的时空称为弯曲时空，以区别于平直时空；在平直时空中，上述变换是可能的.

然而，在非伽利略时空中的任何个别点，我们可以用适当的坐标变换将 $g_{ik}$ 各量化为伽利略坐标的形式：这会导致二次型化为恒定系数（在给定点的 $g_{ik}$ 值）的对角形式，我们把这样的坐标系称为对于一个给定点的伽利略坐标系\footnote{为了避免误解，我们必须立即指出，选取这样的参考系并不意味着在相应的无穷小的四维体积元中消除引力场. 由于等效原理，这样的消除才总有可能实现，且它有着更深刻的涵义（见§87）.}.

我们注意到在一个给定点将 $g_{ik}$ 化为对角形式之后，$g_{ik}$ 各量的矩阵有一个正的主值和三个负的主值（这些正负号的总和称为矩阵的号差）.由此可以断定，由 $g_{ik}$ 各量所构成的行列式 $g$ 在现实的时空中总是负的：
\begin{equation}
  g<0.
\end{equation}

时空度规的变化也就意味着纯粹的空间度规的变化.平直时空中伽利略坐标的 $g_{ik}$ 对应空间欧几里得几何.在引力场中空间的几何变成非欧几里得的.这既适用于时空“弯曲”的“真实”引力场，也适用于时空保持平直、仅仅由于参考系是非惯性参考系而产生的场.

关于引力场中的空间几何问题将会在§84中更详细地考察.这里最好是先进行简单的讨论，直观地说明在改变到非惯性参考系时，空间几何将不可避免地变为非欧几里得几何.考察两个参考系，其中之一（K）是惯性参考系，而另一个（K$'$）相对于 K 围绕着共同的坐标轴 z 匀速旋转.参考系 K 的平面 $xy$ 上的一个圆（其圆心为坐标系的原点）也可以看做参考系 K$'$ 的平面 $x'y'$ 上的一个圆.用一把有刻度的尺子测量这个圆的周长和它的直径，我们将得到比值为 π 的两个数值，这与惯性参考系中的欧几里得几何相符.假设现在的测量过程使用与参考系 K$'$ 相对静止的标尺，并且从参考系 K 中观察这个过程，我们将会看到，沿着圆周放置的标尺会产生洛伦兹收缩，而沿着径向放置的标尺保持不变.因此显而易见，在这样的测量方式下获得的圆周长和直径的比值会大于 π.

在一般情况下任意变化的引力场的空间度规不但是非欧几里得的，而且还随时间变化.这意味着，不同几何距离之间的关系随时间变化.被引入到场中的“试验粒子”的相对位置在任何一个坐标系中都不可能是一成不变的\footnote{严格地说，粒子数应当大于 4. 因为从任何 4 个粒子之间的 6 条线段我们可以构造一个4面体，我们总可以通过适当定义参考系，使 4 个粒子的系统构成一个不变的 4 面体. 何况，在3 个或 2 个粒子的系统中，我们可以使粒子彼此相对固定.}.这样，如果粒子分布在某个圆上和这个圆的直径上，由于圆的周长和直径的比值不等于 π 并且随时间变化，显然，如果粒子沿着直径的距离保持不变，那么沿着圆周的距离应该变化，反之亦然.这样一来，在广义相对论中，一般说来，物体系统不可能是相互静止的.

与狭义相对论中的参考系相比，广义相对论中参考系这个概念本身发生了根本的改变.在狭义相对论中参考系可以理解为相互静止的、相互分布保持不变的物体的总和.在引入变化的引力场的情况下这样的物体系统不存在，为了精确确定粒子在空间中的位置，严格来讲，必须具备充满整个空间的无穷多数量的物体的总和，类似于某种“介质”.这样的物体系统和与每个物体联系在一起的以任意方式运行的时钟构成广义相对论中的参考系.

由于参考系选取的任意性，在广义相对论中对自然规律的描述在形式上应该适用于任意的四维坐标系（即，所谓的“协变”形式）.但是这种状况并不意味着所有的这些参考系在物理上是等价的（类比于狭义相对论中所有惯性参考系的物理等价性）.正相反，具体的物理现象，包括物体运动的性质，在所有的参考系中是不同的.

\section{曲线坐标}

因为在研究引力场时需要考虑在任意参考系中的现象，就必须发展任意曲线坐标中的四维几何.在§83，§85，§86我们将专门讨论这一问题.

让我们研究从一个坐标系 $x^0,x^1,x^2,x^3$ 到另一个坐标系 $x'^0,x'^1,x'^2,x'^3$ 的变换：
\begin{equation*}
  x^i=f^i(x'^0,x'^1,x'^2,x'^3),
\end{equation*}
其中 $f^i$ 是某些函数.在变换坐标时，坐标的微分按照下面的关系式变换：
\begin{equation}
  \mathrm{d}x^i=\frac{\partial x^i}{\partial x'^k}\mathrm{d}x'^k.
\end{equation}

如果任何四个量 $A^i$ 的集合，在坐标变换时像坐标的微分一样地变换，那么，这四个量的集合称为四维逆变矢量：
\begin{equation}
  A^i=\frac{\partial x^i}{\partial x'^k}A'^k.
\end{equation}

设 $\varphi$ 是某一标量.在坐标变换时，四个量 $\partial\varphi/\partial x^i$ 按照公式
\begin{equation}
  \frac{\partial\varphi}{\partial x^i}
  =\frac{\partial\varphi}{\partial x'^k}
  \frac{\partial x'^k}{\partial x^i}
\end{equation}
变换，这个公式与(83.2)式不同.如果四个量的集合 $A_i$ 在坐标变换时像一个标量的导数一样变换，那么，这四个量的集合称为四维协变矢量：
\begin{equation}
  A_i=\frac{\partial x'^k}{\partial x^i}A'_k.
\end{equation}

可以采用类似的方式确定不同阶数的四维张量.于是，像两个逆变矢量的乘积一样变换，也就是按照法则
\begin{equation}
  A^{ik}=\frac{\partial x^i}{\partial x'^l}
  \frac{\partial x^k}{\partial x'^m}A'^{lm}
\end{equation}
变换的16个量的集合称为二阶四维逆变张量 $A^{ik}$.同理，我们定义按下式变换的16个量的集合为二阶协变张量：
\begin{equation}
  A_{ik}=\frac{\partial x'^l}{\partial x^i}
  \frac{\partial x'^m}{\partial x^k}A'_{lm},
\end{equation}
而四维混合张量 $A^i{}_k$ 按照下式定义：
\begin{equation}
  A^i{}_k=\frac{\partial x^i}{\partial x'^l}
  \frac{\partial x'^m}{\partial x^k}A'^l{}_m.
\end{equation}

以上给出的定义是伽利略坐标中的四维矢量和四维张量定义的自然扩充(§6)，按照这些定义，微分 $\mathrm{d}x^i$ 也是逆变的四维矢量，而导数 $\partial\varphi/\partial x^i$ 是协变的四维矢量\footnote{不过，在伽利略坐标系中，坐标 $x^i$ 本身（而不仅仅是其微分）也构成一个四维矢量，在曲线坐标中，情况当然就不是如此了.}.

在曲线坐标中通过将其他四维张量进行互乘或缩并的方式构造四维张量的法则与在伽利略坐标中相同.容易证实，按照变换法则（83.2)和（83.4)，两个四维矢量的标积 $A^iB_i$ 确实是不变的：
\begin{equation*}
  A^iB_i=\frac{\partial x^i}{\partial x'^n}
  \frac{\partial x'^m}{\partial x^i}A^nB'_m
  =\frac{\partial x'^m}{\partial x'^l}A'^lB'_m=A'^lB'_l.
\end{equation*}

在过渡到曲线坐标时单位四维张量 $\delta^i_k$ 的定义仍旧不变：它的分量当 i≠k 时 $\delta^i_k=0$，而当 i=k 时，则等于1.假如 $A^k$ 是一个四维矢量，那么，用 $\delta^i_k$ 乘之，就得到
\begin{equation*}
  A^k\delta^i_k=A^i,
\end{equation*}
它仍然是一个四维矢量；这也证明了 $\delta^i_k$ 是一个张量.

线元的平方 $\mathrm{d}s^2$ 是微分 $\mathrm{d}x^i$ 的二次型，也就是
\begin{equation}
  \mathrm{d}s^2=g_{ik}\mathrm{d}x^i\mathrm{d}x^k,
\end{equation}
其中 $g_{ik}$ 是坐标的函数，$g_{ik}$ 对于指标 i 和 k 是对称的，就是说：
\begin{equation}
  g_{ik}=g_{ki}.
\end{equation}
既然 $g_{ik}$ 与逆变张量 $\mathrm{d}x^i\mathrm{d}x^k$ 的积（缩并）是一个标量，那么，$g_{ik}$ 就是一个协变张量.张量 $g_{ik}$ 称为度规张量.

若两个张量 $A_{ik}$ 和 $B^{ik}$ 满足
\begin{equation*}
  A_{ik}B^{kl}=\delta^l_i,
\end{equation*}
则称为互逆的（互为倒数）.特别指出，张量 $g_{ik}$ 的逆张量 $g^{ik}$ 称为逆变度规张量，即
\begin{equation}
  g_{ik}g^{kl}=\delta^l_i.
\end{equation}

同一个物理量既可以用一组逆变分量来表示，也可以用一组协变分量来表示.显然，唯一能决定逆变分量和协变分量之间关系的一组数值是度规张量的分量.这样的联系由下列公式给出：
\begin{equation}
  A^i=g^{ik}A_k,\qquad A_i=g_{ik}A^k.
\end{equation}

在伽利略坐标系中度规张量具有如下分量：
\begin{equation}
  g^{(0)}_{ik}=g^{ik(0)}=
  \begin{pmatrix}
    1 & 0 & 0 & 0\\
    0 & -1 & 0 & 0\\
    0 & 0 & -1 & 0\\
    0 & 0 & 0 & -1
  \end{pmatrix}.
\end{equation}

公式（83.11）给出已知的关系 $A^0=A_0$，$A^{1,2,3}=-A_{1,2,3}$.\footnote{在我们用伽利略坐标系作类比时，应当认识到只有在平直的四维空间中才能选择这样的坐标系. 对于四维弯曲空间必须给定一个无限小四维体积元，在其内才能言及、且总能找到伽利略坐标系. 所有结论均不受这一改变的影响.}

上述内容也适用于张量.同一物理张量的不同形式之间的转换可以按照下面的公式，借助于度规张量来实现：
\begin{equation*}
  A^i{}_k=g^{il}A_{lk},\qquad
  A^{ik}=g^{il}g^{km}A_{lm},
\end{equation*}
等等.

§6（在伽利略坐标系中）定义了完全反对称的单位赝张量 $e^{iklm}$.我们将其转换到任意的曲线坐标系中，并表示为 $E^{iklm}$.同时仍旧按照以前定义的符号 $e^{iklm}$ 取值 $e^{0123}=1$（或者 $e_{0123}=-1$）.

令 $x'^i$ 为伽利略坐标，而 $x^i$ 为任意的曲线坐标.按照通用的张量变换规则，我们有
\begin{equation*}
  E^{iklm}=\frac{\partial x^i}{\partial x'^p}
  \frac{\partial x^k}{\partial x'^r}
  \frac{\partial x^l}{\partial x'^s}
  \frac{\partial x^m}{\partial x'^t}e^{prst},
\end{equation*}
或者
\begin{equation*}
  E^{iklm}=Je^{iklm},
\end{equation*}
其中 $J$ 为由导数 $\partial x^i/\partial x'^p$ 组成的行列式，它正是从伽利略坐标向曲线坐标变换的雅可比行列式：
\begin{equation*}
  J=\frac{\partial(x^0,x^1,x^2,x^3)}
  {\partial(x'^0,x'^1,x'^2,x'^3)}.
\end{equation*}

这个雅可比行列式可以用由度规张量 $g_{ik}$（在参考系 $x^i$ 中）的行列式来表示.为此我们写出度规张量变换的公式：
\begin{equation*}
  g^{ik}=\frac{\partial x^i}{\partial x'^l}
  \frac{\partial x^k}{\partial x'^m}g^{lm(0)},
\end{equation*}
并且令这个等式两边的值所构成的行列式也相等.逆张量的行列式 $|g^{ik}|=1/g$.行列式 $|g^{lm(0)}|=-1$.因此有 $1/g=-J^2$，由此 $J=1/\sqrt{-g}$.

这样一来，曲线坐标系中四阶反对称的单位张量必须定义为
\begin{equation}
  E^{iklm}=\frac{1}{\sqrt{-g}}e^{iklm}.
\end{equation}

可以利用下面的公式将该张量的指标下移：
\begin{equation*}
  e^{prst}g_{ip}g_{kr}g_{ls}g_{mt}=-ge_{iklm},
\end{equation*}
这样它的协变分量是
\begin{equation}
  E_{iklm}=\sqrt{-g}\,e_{iklm}.
\end{equation}

在伽利略坐标系 $x'^i$ 中，一个标量对于 $\mathrm{d}\varOmega'=\mathrm{d}x'^0\mathrm{d}x'^1\mathrm{d}x'^2\mathrm{d}x'^3$ 的积分也是标量.就是说单元 $\mathrm{d}\varOmega'$ 在积分中具有标量的性质（§6）.在向曲线坐标 $x^i$ 变换时积分单元 $\mathrm{d}\varOmega'$ 变为
\begin{equation*}
  \mathrm{d}\varOmega'\to\frac{1}{J}\mathrm{d}\varOmega
  =\sqrt{-g}\,\mathrm{d}\varOmega.
\end{equation*}

因此，在曲线坐标中，在四维空间的某一区域内积分时，$\sqrt{-g}\,\mathrm{d}\varOmega$ 具有不变量的特征\footnote{如果 $\varphi$ 是一个标量，则量 $\sqrt{-g}\,\varphi$ 对 $\mathrm{d}\varOmega$ 积分构成一不变量，称为标量密度. 类似地，我们称 $\sqrt{-g}\,A^i$ 为矢量密度，$\sqrt{-g}\,A^{ik}$ 为张量密度，等等. 这些量在乘以四维体积元 $\mathrm{d}\varOmega$ 后成为矢量或张量（一般说来，沿有限体积的积分 $\int A^i\sqrt{-g}\,\mathrm{d}\varOmega$ 不可能是矢量，因为矢量 $A^i$ 的变换法则在不同的点是不同的）.}.

在§6末关于在超曲面上的积分元、曲面上的积分元、曲线上的积分元等的论述对于曲线坐标也保持有效，只有一点例外，即对偶张量的定义有些改变.由三个无限小线段所构成的超曲面的“面积”元是一个逆变反对称张量 $\mathrm{d}S^{ikl}$；与其对偶的矢量可用张量 $\sqrt{-g}\,e_{iklm}$ 乘之得到，即等于
\begin{equation}
  \sqrt{-g}\,\mathrm{d}S_i
  =-\frac{1}{6}e_{iklm}\mathrm{d}S^{klm}\sqrt{-g}.
\end{equation}

同理，假如 $\mathrm{d}f^{ik}$ 是由两个无穷小的线段张成的（两维）面元，那么，与其对偶的张量定义为\footnote{意思是，不论坐标 $x^i$ 的几何意义怎样，由无穷小位移 $\mathrm{d}x^i$, $\mathrm{d}x'^{i}$, $\mathrm{d}x''^{i}$ 构造面元 $\mathrm{d}S^{ikl}$ 和 $\mathrm{d}f^{ik}$ 的方式与§6中是相同的. 那么，$\mathrm{d}S_i$ 和 $\mathrm{d}f^{*}_{ik}$ 的形式意义就同以前一样了. 特别是，同以前一样，$\mathrm{d}S_0=\mathrm{d}x^1\mathrm{d}x^2\mathrm{d}x^3=\mathrm{d}V$. 我们保持 $\mathrm{d}V$ 为 3 个空间坐标微分乘积的较早期定义；然而必须记住，在曲线坐标中，几何上的空间体积元不是由 $\mathrm{d}V$ 而是由 $\sqrt{\gamma}\,\mathrm{d}V$ 给出，这里 $\gamma$ 是空间度规张量的行列式（将在下节给出）.}
\begin{equation}
  \sqrt{-g}\,\mathrm{d}f^{*}_{ik}
  =\frac{1}{2}\sqrt{-g}\,e_{iklm}\mathrm{d}f^{lm}.
\end{equation}

我们像以前一样在此用 $\mathrm{d}S_i$ 和 $\mathrm{d}f^{*}_{ik}$ 分别代表 $\frac{1}{6}e_{iklm}\mathrm{d}S^{klm}$ 和 $\frac{1}{2}e_{iklm}\mathrm{d}f^{lm}$（不是它们与 $\sqrt{-g}$ 的乘积）；各种积分彼此变换的法则(6.14)—(6.19)也保持不变，因为它们只是一些形式符号的推导，而与相应的量的张量性质无关.其中，有一个变换法则是特别重要的，这就是将在一个超曲面上的积分变换为在一个四维体积上的积分的变换法则（高斯定理），这个变换可以用下面的替换来实现：
\begin{equation}
  \mathrm{d}S_i\to\mathrm{d}\varOmega\frac{\partial}{\partial x^i}.
\end{equation}

\section{距离与时间间隔}

我们已经说过，在广义相对论中，对于坐标系的选择并未加任何限制；三个空间坐标 $x^1,x^2,x^3$，可以是决定物体在空间中的位置的任何量，而时间坐标 $x^0$ 则可以用一个任意行走的钟来确定.因此就有这样一个问题发生，即如何用 $x^0,x^1,x^2,x^3$ 这些量的值来表示实在的距离和时间间隔.

首先，我们找出固有时（以后我们用 $\tau$ 来表示它）与坐标 $x^0$ 的关系.为此，我们来考虑在空间的同一点发生的两个无限近的事件.如我们所知道的，两个事件的间隔 ds 恰恰是 $\mathrm{d}s=c\,\mathrm{d}\tau$，此处的 $\mathrm{d}\tau$ 是两个事件之间的时间（固有时)间隔.因此，在普遍式子 $\mathrm{d}s^2=g_{ik}\mathrm{d}x^i\mathrm{d}x^k$ 内，设 $\mathrm{d}x^1=\mathrm{d}x^2=\mathrm{d}x^3=0$ 我们便求得
\begin{equation*}
  \mathrm{d}s^2=c^2\mathrm{d}\tau^2
  =g_{00}(\mathrm{d}x^0)^2,
\end{equation*}
从而
\begin{equation}
  \mathrm{d}\tau=\frac{1}{c}\sqrt{g_{00}}\,\mathrm{d}x^0,
\end{equation}
或者，在空间的同一点发生的任意两个事件之间的时间
\begin{equation}
  \tau=\frac{1}{c}\int\sqrt{g_{00}}\,\mathrm{d}x^0.
\end{equation}

这个关系式决定相应于坐标 $x^0$ 的变化的实际时间间隔（或称对于空间一给定点的固有时).我们注意到，从这些公式可以看出，量 $g_{00}$ 是正的，即
\begin{equation}
  g_{00}>0.
\end{equation}

必须强调指出下面两个条件的意义的差别，第一个条件是(84.3)，第二个条件是关于由张量 $g_{ik}$ 三个主值的正负号决定的号差（§82).不满足第二个条件的张量 $g_{ik}$ 一般不能与任何真实的引力场对应，也就是说它不能是一个真实的时空的度规.不满足条件(84.3)只表明相应的参考系不能用真实的物体来实现，假如这时加在主值上的条件得到满足，那么一个适当的坐标变换就可以使 $g_{00}$ 变为正（旋转坐标系就是这种参考系的一个例子，见§89).

现在我们来求空间距离元 dl.在狭义相对论中，我们可定义 dl 为同时发生、相距无限近的两个事件之间的间隔.在广义相对论中，这在通常的情况下是不可能的，也就是说，简单地使在 ds 中的 $\mathrm{d}x^0=0$ 来决定 dl 是不可能的.这是因为在引力场中，对于空间的不同点，固有时与坐标 $x^0$ 有不同的关系.

为了求 dl，现在我们的做法如下.

假设一个光信号从空间的一给定点 B（具有坐标 $x^{\alpha}+\mathrm{d}x^{\alpha}$)跑向相距无限近的另一点 A（具有坐标 $x^{\alpha}$)，然后又沿原路回去.显然（从空间点 B 观察）为此所必需的时间乘以 c 是两点间的距离的两倍.

我们写出间隔，并将时间坐标和空间坐标分开：
\begin{equation}
  \mathrm{d}s^2=g_{\alpha\beta}\mathrm{d}x^{\alpha}\mathrm{d}x^{\beta}
  +2g_{0\alpha}\mathrm{d}x^0\mathrm{d}x^{\alpha}
  +g_{00}(\mathrm{d}x^0)^2,
\end{equation}
不言而喻，此处任何重复两次的希腊字母指标，我们都从1到3求和.对应信号从一点出发和该信号到达另一点的两个事件的间隔为零.相对于 $\mathrm{d}x^0$ 求解方程 $\mathrm{d}s^2=0$，我们求得两个根：
\begin{equation}
\begin{aligned}
  \mathrm{d}x^{0(1)}
  &=\frac{1}{g_{00}}\left(-g_{0\alpha}\mathrm{d}x^{\alpha}
  -\sqrt{(g_{0\alpha}g_{0\beta}-g_{\alpha\beta}g_{00})
  \mathrm{d}x^{\alpha}\mathrm{d}x^{\beta}}\right),\\
  \mathrm{d}x^{0(2)}
  &=\frac{1}{g_{00}}\left(-g_{0\alpha}\mathrm{d}x^{\alpha}
  +\sqrt{(g_{0\alpha}g_{0\beta}-g_{\alpha\beta}g_{00})
  \mathrm{d}x^{\alpha}\mathrm{d}x^{\beta}}\right),
\end{aligned}
\end{equation}
它们分别对应 A 和 B 之间两个方向上的信号传播.如果 $x^0$ 为信号到达 A 的时刻，那么它从 B 点出发和回到 B 点的时刻分别为 $x^0+\mathrm{d}x^{0(1)}$ 和 $x^0+\mathrm{d}x^{0(2)}$.在示意图18中实线为给定的坐标 $x^{\alpha}$ 和 $x^{\alpha}+\mathrm{d}x^{\alpha}$ 所对应的世界线，而虚线是信号的世界线\footnote{在图18上假设了 $\mathrm{d}x^{0(2)}>0$, $\mathrm{d}x^{0(1)}<0$，然而这不是必须的：$\mathrm{d}x^{0(1)}$ 和 $\mathrm{d}x^{0(2)}$ 可以具有同一个符号. 事实在于，在这种条件下，在信号到达 $A$ 时刻 $x^{0}(A)$ 的值可以小于信号从 $B$ 出发时刻 $x^{0}(B)$ 的值，这不会造成任何的自相矛盾，因为在不同空间点内的时钟的进程并没有以某种方式设定为同步的.}.因此，信号从某一点出发而又回到该点的总“时间”间隔等于
\begin{equation*}
  \mathrm{d}x^{0(2)}-\mathrm{d}x^{0(1)}
  =\frac{2}{g_{00}}\sqrt{(g_{0\alpha}g_{0\beta}
  -g_{\alpha\beta}g_{00})\mathrm{d}x^{\alpha}\mathrm{d}x^{\beta}}.
\end{equation*}

\par\nopagebreak\noindent\begin{minipage}{\linewidth}\centering\includegraphics[width=4cm]{fig18.pdf}\\[0.3em]{\small 图 18}\end{minipage}\par\nopagebreak

按照(84.1)，将上式乘以 $\sqrt{g_{00}}/c$，就得到相应的固有时的间隔，再乘以 $c/2$，就得到两点间的距离 dl.我们得到的结果是
\begin{equation*}
  \mathrm{d}l^2=\left(-g_{\alpha\beta}
  +\frac{g_{0\alpha}g_{0\beta}}{g_{00}}\right)
  \mathrm{d}x^{\alpha}\mathrm{d}x^{\beta}.
\end{equation*}
这就是所要求的用空间坐标元来确定距离的式子.我们把它改写如下：
\begin{equation}
  \mathrm{d}l^2=\gamma_{\alpha\beta}\mathrm{d}x^{\alpha}\mathrm{d}x^{\beta},
\end{equation}
其中，
\begin{equation}
  \gamma_{\alpha\beta}=-g_{\alpha\beta}
  +\frac{g_{0\alpha}g_{0\beta}}{g_{00}}
\end{equation}
是三维度规张量，它决定空间的度规，亦即决定空间的几何性质.关系式 (84.7) 联系现实的空间的度规和四维时空的度规\footnote{二次型（84.6）应该必须是正的. 为此它的系数应该像已知的那样，满足条件
\begin{equation*}
  \gamma_{11}>0,\quad
  \begin{vmatrix}\gamma_{11} & \gamma_{12}\\ \gamma_{21} & \gamma_{22}\end{vmatrix}>0,\quad
  \begin{vmatrix}\gamma_{11} & \gamma_{12} & \gamma_{13}\\
  \gamma_{21} & \gamma_{22} & \gamma_{23}\\
  \gamma_{31} & \gamma_{32} & \gamma_{33}\end{vmatrix}>0.
\end{equation*}
用 $g_{ik}$ 表示 $\gamma_{ik}$，容易求得，这些条件具有如下形式
\begin{equation*}
  \begin{vmatrix}g_{00} & g_{01}\\ g_{10} & g_{11}\end{vmatrix}<0,\quad
  \begin{vmatrix}g_{00} & g_{01} & g_{02}\\ g_{10} & g_{11} & g_{12}\\
  g_{20} & g_{21} & g_{22}\end{vmatrix}>0,\quad g<0.
\end{equation*}
在整个参考系内，度规张量的分量应该满足这些条件以及条件(84.3)，这个参考系可以用实在的物体来实现.}.

然而，我们必须记着，$g_{ik}$ 一般与 $x^0$ 有关，因而，空间度规（84.7）也随着时间变化.根据这个理由，将 dl 积分就没有意义；这个积分与空间中给定的两点间的世界线（沿着它积分）有关.因此，一般来说，在广义相对论中，物体间的一定距离的概念失掉了意义，仅仅对于无限小的距离还保持有效.只有在 $g_{ik}$ 与时间无关的参考系中，才能在空间的一个有限区域内确定距离，这时沿着一条空间曲线的积分 $\int\mathrm{d}l$ 才有一定的意义.

有必要指出，张量 $-\gamma_{\alpha\beta}$ 是三维逆变张量 $g^{\alpha\beta}$ 的逆张量.实际上，将等式 $g^{ik}g_{kl}=\delta^i_l$ 以分量的形式展开，我们有：
\begin{equation}
\begin{aligned}
  &g^{\alpha\beta}g_{\beta\gamma}+g^{\alpha 0}g_{0\gamma}=\delta^{\alpha}_{\gamma},\\
  &g^{\alpha\beta}g_{\beta 0}+g^{\alpha 0}g_{00}=0,\\
  &g^{0\beta}g_{\beta 0}+g^{00}g_{00}=1.
\end{aligned}
\end{equation}
从第二个等式确定 $g^{\alpha 0}$ 然后代入第一个等式，得到：
\begin{equation*}
  -g^{\alpha\beta}\gamma_{\beta\gamma}=\delta^{\alpha}_{\gamma},
\end{equation*}
这就是需要证明的结果.这个结果还可以用其他形式表述为，量 $-g^{\alpha\beta}$ 构成对应度规（84.7）的三维逆变度规张量：
\begin{equation}
  \gamma^{\alpha\beta}=-g^{\alpha\beta}.
\end{equation}

我们同样指出，由量 $g_{ik}$ 和 $\gamma_{\alpha\beta}$ 分别构成的行列式 $g$ 和 $\gamma$ 相互之间可由简单的关系式联系起来：
\begin{equation}
  -g=g_{00}\gamma.
\end{equation}

在接下来的一系列应用中，引入三维矢量 $\boldsymbol{g}$ 会带来一定的方便.矢量 $\boldsymbol{g}$ 的协变分量由下式确定：
\begin{equation}
  g_{\alpha}=-\frac{g_{0\alpha}}{g_{00}}.
\end{equation}

将 $\boldsymbol{g}$ 看做带有度规(84.7)的空间中的矢量，我们必须将它的逆变分量定义为 $g^{\alpha}=\gamma^{\alpha\beta}g_{\beta}$.基于（84.9)和（84.8)的第二个等式容易看出，
\begin{equation}
  g^{\alpha}=\gamma^{\alpha\beta}g_{\beta}=-g^{0\alpha}.
\end{equation}
从（84.8)的第三个等式同样可以得到公式
\begin{equation}
  g^{00}=\frac{1}{g_{00}}-g_{\alpha}g^{\alpha}.
\end{equation}

现在我们再来讨论广义相对论中“同时”概念的定义，换句话说，我们讨论处于空间不同点的钟能否同步的问题，也就是这些钟的对时问题.

显然，这样的同步应该通过两点之间光信号的交换来实现.我们再来考察两个无限接近的点 A 和 B 之间光信号的传播过程，如图18中所示.位于 B 点的时钟记录了信号由 B 点出发时刻的读数和信号返回 B 点时刻的读数，我们必须认为这两个读数的中间值与信号到达 A 点的时刻 $x^0$ 是同时的，该时刻为
\begin{equation*}
  x^0+\Delta x^0=x^0+\frac{1}{2}\left(\mathrm{d}x^{0(2)}
  +\mathrm{d}x^{0(1)}\right).
\end{equation*}
将（84.5）代入上式，找到发生在两个无限接近的点的两个同时事件的“时间”$x^0$ 值的差，具有下列形式，
\begin{equation}
  \Delta x^0=-\frac{g_{0\alpha}\mathrm{d}x^{\alpha}}{g_{00}}
  \equiv g_{\alpha}\mathrm{d}x^{\alpha}.
\end{equation}

这个关系使我们能够在空间的任何无限小的体积内来校准钟.从 A 点继续作类似的校准，我们可以沿任何（非闭合的）曲线校准不同的钟，也就是说我们能够确定事件的同时性\footnote{给等式（84.14）乘以 $g_{00}$ 并且将两个项移到等式的一边，可以将同步条件表示为 $\mathrm{d}x_0=g_{0i}\mathrm{d}x^i=0$；两个无限接近的同时发生的事件之间的“协变微分”$\mathrm{d}x_0$ 应该等于零.}.

然而，一般来说，沿着某条闭合路径来校准时间是无法做到的.事实上，沿着闭合路径行进，又回到原来出发的那一点，我们会得到一个不等于零的 $\Delta x^0$ 值.要在整个空间中统一地同步时间尤其是不可能的，只有在 $g_{0\alpha}$ 各量都为零的参考系中才是例外\footnote{这里还应该补充一种情况，就是当 $g_{0\alpha}$ 可以通过时间坐标的简单变换化为零，并且这种变换不能触及用于确定空间坐标的物体系统的选取.}.

应该强调，不可能对所有的时钟同步是任意参考系的特征，而不是时空本身的特征.在任意一个引力场中总是可以选择（甚至无穷多的方法）这样的参考系，使得 $g_{0\alpha}$ 的三个值恒等于0，进而使得时钟的完全同步成为可能(见§97).

在狭义相对论中，两个相对运动的时钟的固有时流动已经不同.在广义相对论中，同一参考系里的不同空间点的固有时以不同的方式流动.这意味着，发生在某一空间点的两个事件之间的固有时间隔，与发生在另一空间点，并与前两个事件同时发生的另两个事件之间的时间间隔，一般说来也是不同的.

\section{协变微分}

在伽利略坐标系中\footnote{更一般地，只要 $g_{ik}$ 各量是常数.}，矢量 $A_i$ 的微分 $\mathrm{d}A_i$ 仍然是矢量，而矢量的分量对坐标的导数 $\partial A_i/\partial x^k$ 构成一个张量.在曲线坐标中就不是这样；$\mathrm{d}A_i$ 不是矢量，$\partial A_i/\partial x^k$ 不是张量，这是由于 $\mathrm{d}A_i$ 是处在空间不同点（相距无限近）上的矢量之差，而在这些空间不同点上，矢量的变换是不同的，因为变换公式(83.2),(83.4) 中的系数是坐标的函数.

这些都不难直接证明.为此，我们求微分 $\mathrm{d}A_i$ 在曲线坐标中的变换公式.一个协变矢量是按照公式
\begin{equation*}
  A_i=\frac{\partial x'^k}{\partial x^i}A'_k
\end{equation*}
变换的，因而
\begin{equation*}
  \mathrm{d}A_i=\frac{\partial x'^k}{\partial x^i}\mathrm{d}A'_k
  +A'_k\mathrm{d}\frac{\partial x'^k}{\partial x^i}
  =\frac{\partial x'^k}{\partial x^i}\mathrm{d}A'_k
  +A'_k\frac{\partial^2 x'^k}{\partial x^i\partial x^l}\mathrm{d}x^l.
\end{equation*}
因此，$\mathrm{d}A_i$ 根本不像矢量一样变换(逆变矢量的微分当然也是这样).仅仅当二阶导数 $\partial^2 x'^k/\partial x^i\partial x^l=0$ 时，即当 $x'^k$ 是 $x^k$ 的线性函数时，变换公式有以下的形式：
\begin{equation*}
  \mathrm{d}A_i=\frac{\partial x'^k}{\partial x^i}\mathrm{d}A'_k,
\end{equation*}
就是说，$\mathrm{d}A_i$ 像矢量一样变换.

现在我们研究张量的定义，其在曲线坐标内所起的作用，正如 $\partial A_i/\partial x^k$ 在伽利略坐标中所起的作用一样，换句话说，我们应当将 $\partial A_i/\partial x^k$ 从伽利略坐标改换到曲线坐标中去.

在曲线坐标中，为了得到矢量的微分，而这个微分要有矢量的特性，那么，两个相减的矢量就必须在空间的同一点.换句话说，我们必须用某一个方法将彼此相距无限近的两个矢量中的一个“移”到第二个矢量所在之点，在此之后，我们求两个矢量的差，现在这两个矢量是处在空间的同一点.这个移动运算的本身必须如此定义，使这个差在伽利略坐标中与普通微分 $\mathrm{d}A_i$ 相合.既然 $\mathrm{d}A_i$ 不过是两个相距无限近的矢量的分量之差，这就意味着，当我们用伽利略坐标时，矢量的分量经过移动的操作应当不发生改变.而这样的移动正是矢量与它自己平行的移动.当矢量作平行移动时，它在伽利略坐标中的分量不变；假如应用曲线坐标，那么，一般来说，在这样移动时，矢量的分量要改变.因此，在曲线坐标中，在将一个矢量平行移动到另一点以后，两个矢量的分量之差与移动前两者之差(即微分 $\mathrm{d}A_i$)不相等.

所以，为了比较两个相距无限近的矢量，我们应当将其中之一平行移动到第二个矢量所在的点.让我们考虑一个任意的逆变矢量；假如它在 $x^i$ 点的值是 $A^i$，那么，在邻近之点 $x^i+\mathrm{d}x^i$，它的值就是 $A^i+\mathrm{d}A^i$.我们将矢量 $A^i$ 作无限小的平行移动，移到 $x^i+\mathrm{d}x^i$ 点，这时矢量的变化用 $\delta A^i$ 来表示.现在位于同一点的两个矢量之差用 $\mathrm{D}A^i$ 来表示，则 $\mathrm{D}A^i$ 等于
\begin{equation}
  \mathrm{D}A^i=\mathrm{d}A^i-\delta A^i.
\end{equation}

一个矢量在作无限小的平行移动时，其分量的变化 $\delta A^i$ 与分量本身的值有关，且这个关系显然应当是线性的.这点可以直接从下述事实来推断，即两个矢量之和必须按照每个矢量的变换规则来变换.因此，$\delta A^i$ 有下面的形式：
\begin{equation}
  \delta A^i=-\Gamma^i_{kl}A^k\mathrm{d}x^l,
\end{equation}
其中 $\Gamma^i_{kl}$ 是坐标的函数，其形式当然与坐标系有关，在伽利略坐标系中，$\Gamma^i_{kl}=0$.

由此已可看出，$\Gamma^i_{kl}$ 各量并不构成一个张量，因为假如在一个坐标系中一个张量为零，那么，在任何其他坐标系中它都为零.在曲线坐标中当然不可能使所有 $\Gamma^i_{kl}$ 在全空间中都为零.

但是等效原理要求通过适当的坐标系选取能够在给定的无限小空间区域消除引力场，也就是说使 $\Gamma^i_{kl}$ 各量为零.在下面§87节我们会看到 $\Gamma^i_{kl}$ 扮演了场强的角色\footnote{当我们谈及伽利略坐标系时，所指的其实是这样的坐标系，为了简便起见才简单地称之为“伽利略坐标系”. 进一步而言，所有的证明不仅对平直空间成立，对弯曲的四维空间同样成立.}.

$\Gamma^i_{kl}$ 各量称为联络系数或者克里斯托夫符号.

以后我们还要使用 $\Gamma_{i,kl}$ 各量\footnote{有时候用符号 $\left\{\begin{matrix}kl\\ i\end{matrix}\right\}$ 和 $\left[\begin{matrix}kl\\ i\end{matrix}\right]$ 相应地来代替 $\Gamma^i_{kl}$ 和 $\Gamma_{i,kl}$.}，其定义如下：
\begin{equation}
  \Gamma_{i,kl}=g_{im}\Gamma^m_{kl}.
\end{equation}
反过来，
\begin{equation}
  \Gamma^i_{kl}=g^{im}\Gamma_{m,kl}.
\end{equation}

不难建立协变矢量的分量在平行移动下的变化与克里斯托夫符号间的关系.为此，我们指出，一个标量在平行移动时是不变的.就特例而言，两个矢量的标积在平行移动时是不变的.

设 $A_i$ 和 $B^i$ 是任意的协变和逆变矢量，则从 $\delta(A_iB^i)=0$，我们有
\begin{equation*}
  B^i\delta A_i=-A_i\delta B^i
  =\Gamma^i_{kl}B^kA_i\mathrm{d}x^l,
\end{equation*}
交换指标可得
\begin{equation*}
  B^i\delta A_i=\Gamma^k_{il}A_kB^i\mathrm{d}x^l.
\end{equation*}
鉴于 $B^i$ 的任意性，由此可得
\begin{equation}
  \delta A_i=\Gamma^k_{il}A_k\mathrm{d}x^l,
\end{equation}
此式决定协变矢量在平行移动后的改变.

将（85.2）和 $\mathrm{d}A^i=\frac{\partial A^i}{\partial x^l}\mathrm{d}x^l$ 代入（85.1）式可得到
\begin{equation}
  \mathrm{D}A^i=\left(\frac{\partial A^i}{\partial x^l}
  +\Gamma^i_{kl}A^k\right)\mathrm{d}x^l.
\end{equation}
同理，对于一个协变矢量，可求得
\begin{equation}
  \mathrm{D}A_i=\left(\frac{\partial A_i}{\partial x^l}
  -\Gamma^k_{il}A_k\right)\mathrm{d}x^l.
\end{equation}

(85.6)和（85.7)两式的括号内的表达式是张量，因为乘以矢量 $\mathrm{d}x^l$ 后，其结果是矢量.显然，这些张量就是导数概念在曲线坐标中的推广.这些张量分别称为矢量 $A^i$ 和 $A_i$ 的协变导数.我们用 $A^i{}_{;k}$ 和 $A_{i;k}$ 来代表它们.于是，
\begin{equation}
  \mathrm{D}A^i=A^i{}_{;l}\mathrm{d}x^l,\qquad
  \mathrm{D}A_i=A_{i;l}\mathrm{d}x^l,
\end{equation}
而协变导数本身是
\begin{equation}
  A^i{}_{;l}=\frac{\partial A^i}{\partial x^l}+\Gamma^i_{kl}A^k,
\end{equation}
\begin{equation}
  A_{i;l}=\frac{\partial A_i}{\partial x^l}-\Gamma^k_{il}A_k.
\end{equation}

在伽利略坐标系中，$\Gamma^i_{kl}=0$，协变微分化为普通微分.

计算一个张量的协变导数并不困难.为此，我们必须决定这个张量在无限小的平行移动下的变化.例如，我们来考虑任意一个逆变张量，这个张量是两个逆变矢量之积 $A^iB^k$.在平行移动下，我们有
\begin{equation*}
  \delta(A^iB^k)=A^i\delta B^k+B^k\delta A^i
  =-A^i\Gamma^k_{lm}B^l\mathrm{d}x^m
  -B^k\Gamma^i_{lm}A^l\mathrm{d}x^m.
\end{equation*}
由于这个变换是线性的，对于任意的张量 $A^{ik}$，我们也应当有：
\begin{equation}
  \delta A^{ik}=-\left(\Gamma^k_{ml}A^{im}
  +\Gamma^i_{ml}A^{mk}\right)\mathrm{d}x^l.
\end{equation}
将此式代入
\begin{equation*}
  \mathrm{D}A^{ik}=\mathrm{d}A^{ik}-\delta A^{ik}
  \equiv A^{ik}{}_{;l}\mathrm{d}x^l,
\end{equation*}
我们求得张量 $A^{ik}$ 的协变导数如下：
\begin{equation}
  A^{ik}{}_{;l}=\frac{\partial A^{ik}}{\partial x^l}
  +\Gamma^i_{ml}A^{mk}+\Gamma^k_{ml}A^{im}.
\end{equation}

用完全相似的方法，我们得到混合张量和协变张量的协变导数如下：
\begin{equation}
  A^i_{k;l}=\frac{\partial A^i_k}{\partial x^l}
  -\Gamma^m_{kl}A^i_m+\Gamma^i_{ml}A^m_k,
\end{equation}
\begin{equation}
  A_{ik;l}=\frac{\partial A_{ik}}{\partial x^l}
  -\Gamma^m_{il}A_{mk}-\Gamma^m_{kl}A_{im}.
\end{equation}

用类似的方法，可以决定一个任意阶张量的协变导数.这时，我们得到如下的法则：为了得到张量 $A^{\cdots}_{\cdots}$ 对 $x^l$ 的协变导数，我们在写下普通导数 $\partial A^{\cdots}_{\cdots}/\partial x^l$ 之后，对每个协变指标 $i$（$A^{\cdots}_{.i.}$）加上 $-\Gamma^k_{il}A^{\cdots}_{.k.}$ 一项，而对于每一逆变指标 $i$（$A^{.i.}_{\cdots}$）则加上 $+\Gamma^i_{kl}A^{.k.}_{\cdots}$ 一项.

很容易证明，乘积的协变导数的求法与乘积的普通导数的求法是一样的.为此，我们必须将标量 $\varphi$ 的协变导数看做为普通导数，也就是看做为协变矢量 $\varphi_k=\partial\varphi/\partial x^k$，这和以下事实相符：对于一个标量，$\delta\varphi=0$，因此 $\mathrm{D}\varphi=\mathrm{d}\varphi$.例如，乘积 $A_iB_k$ 的协变导数是
\begin{equation*}
  (A_iB_k)_{;l}=A_{i;l}B_k+A_iB_{k;l}.
\end{equation*}

假如在协变导数中将表示微分的指标上移，我们就得到所谓逆变导数.因此，
\begin{equation*}
  A_i{}^{;k}=g^{kl}A_{i;l},\qquad
  A^{i;k}=g^{kl}A^i{}_{;l}.
\end{equation*}

现在我们来写出克里斯托夫符号从一个坐标系变换到另一个坐标系的变换公式.

比较定义协变导数的方程两边的变换规律，并且要求这些规律对于两边都是一样的，就可以得到这些公式.因此，经简单的计算后可得到
\begin{equation}
  \Gamma^i_{kl}=\Gamma'^m_{np}
  \frac{\partial x^i}{\partial x'^m}
  \frac{\partial x'^n}{\partial x^k}
  \frac{\partial x'^p}{\partial x^l}
  +\frac{\partial^2 x'^m}{\partial x^k\partial x^l}
  \frac{\partial x^i}{\partial x'^m}.
\end{equation}

从这个公式很明显地可以看出，$\Gamma^i_{kl}$ 仅在线性坐标变换下和张量一样变换（这时（85.15）式中的第二项为零）.

但我们发现，上式第二项按照指标 k,l 是对称的，因此在对差值 $S^i_{kl}=\Gamma^i_{kl}-\Gamma^i_{lk}$ 进行坐标变换时就消掉了.因而该差值按照张量法则变换：
\begin{equation*}
  S^i_{kl}=S'^m_{np}\frac{\partial x^i}{\partial x'^m}
  \frac{\partial x'^n}{\partial x^k}
  \frac{\partial x'^p}{\partial x^l},
\end{equation*}
也就是说它是张量，称做空间的挠率张量.

现在我们来证明，在基于等效原理的理论中，挠率张量应该为零.实际上，正如已经说过的，按照等效原理应该存在“伽利略”坐标系，在其中给定的点 $\Gamma^i_{kl}$ 的值为零，相应的 $S^i_{kl}$ 也一样.由于 $S^i_{kl}$ 是张量，那么，在一个坐标系中为零，它将在所有的坐标系中都为零，这意味着克里斯托夫符号按照下指标应该是对称的：
\begin{equation}
  \Gamma^i_{kl}=\Gamma^i_{lk}.
\end{equation}
显然
\begin{equation}
  \Gamma_{i,kl}=\Gamma_{i,lk}.
\end{equation}

在通常情况下总共有40个不同的 $\Gamma^i_{kl}$ 值——对应每一个指标 i，指标 k 和 l 的取值有10对不同的组合（指标 l 和 k 互换位置的一对克里斯托夫符号视做等同）而 i 可以有4种取值.

在(85.16）的条件下公式（85.15)使得我们可以证明上文所作的关于坐标系选择的结论，就是说针对任何一个预先指定的点，总是能选择一个坐标系，使得所有 $\Gamma^i_{kl}$ 的值在该点为零（这样的坐标系称为局部惯性系或局部测地系，见§87\footnote{同时可以指出，通过适当地选取坐标系可以不仅在给定的点上，而且能够沿着给定的世界线将所有的 $\Gamma^i_{kl}$ 化为零（这个论点的证明可以在这本书中找到：P. K. Rashevskii, \textit{Riemannian Geometry and Tensor Analysis}, Nauka（П. К. Рашевский, Риманова геометрия и тензорный анализ）, 1964, §91）.}).

实际上，假设选取坐标原点为指定的点，并且在该点 $\Gamma^i_{kl}$ 具有（在坐标系 $x^i$ 中）初值 $(\Gamma^i_{kl})_0$.在该点的邻域作变换
\begin{equation}
  x'^i=x^i+\frac{1}{2}\left(\Gamma^i_{kl}\right)_0x^kx^l.
\end{equation}
那么
\begin{equation}
  \left(\frac{\partial^2 x'^m}{\partial x^k\partial x^l}
  \frac{\partial x^i}{\partial x'^m}\right)_0
  =\left(\Gamma^i_{kl}\right)_0,
\end{equation}
按照（85.15），所有 $\Gamma'^m_{np}$ 为零.

我们着重强调，在这里条件（85.16）是非常重要的：等式（85.19)左边的表达式按照指标 k,l 是对称的，因此等式的右边部分也应该是对称的.

我们指出，对于变换(85.18)，满足
\begin{equation*}
  \left(\frac{\partial x'^i}{\partial x^k}\right)_0=\delta^i_k,
\end{equation*}
因此它不会改变指定点上的任何一个张量的值（也包括张量 $g_{ik}$），所以克里斯托夫符号的归零和把 $g_{ik}$ 化为伽利略形式可以同时实现.

\section{克里斯托夫符号与度规张量的关系}

我们来证明度规张量 $g_{ik}$ 的协变导数等于零.为此，我们注意到，关系式
\begin{equation*}
  \mathrm{D}A_i=g_{ik}\mathrm{D}A^k
\end{equation*}
对于矢量 $\mathrm{D}A_i$，和对于任何矢量一样，是成立的.另一方面，$A_i=g_{ik}A^k$，所以
\begin{equation*}
  \mathrm{D}A_i=\mathrm{D}(g_{ik}A^k)
  =g_{ik}\mathrm{D}A^k+A^k\mathrm{D}g_{ik}.
\end{equation*}
同 $\mathrm{D}A_i=g_{ik}\mathrm{D}A^k$ 比较，并注意到矢量 $A^i$ 是任意的，就可得到
\begin{equation*}
  \mathrm{D}g_{ik}=0.
\end{equation*}
由此直接推出，协变导数
\begin{equation}
  g_{ik;l}=0.
\end{equation}

因此，在协变微分时，$g_{ik}$ 可以当做常量.

要通过度规张量 $g_{ik}$ 来表示克里斯托夫符号 $\Gamma^i_{kl}$，我们可以利用方程 $g_{ik;l}=0$.为此，按照张量的协变导数的普遍定义(85.14)，我们写出
\begin{equation*}
  g_{ik;l}=\frac{\partial g_{ik}}{\partial x^l}
  -g_{mk}\Gamma^m_{il}-g_{im}\Gamma^m_{kl}
  =\frac{\partial g_{ik}}{\partial x^l}
  -\Gamma_{k,il}-\Gamma_{i,kl}=0.
\end{equation*}
由此，$g_{ik}$ 的导数可用克里斯托夫符号表示出来\footnote{因此，选择局部测地坐标系意味着，在给定点度规张量分量的所有一阶导数均为零.}.我们写出 $g_{ik}$ 的导数值，并将指标 i,k,l 换位，得到：
\begin{equation*}
  \frac{\partial g_{ik}}{\partial x^l}=\Gamma_{k,il}+\Gamma_{i,kl},
  \quad
  \frac{\partial g_{li}}{\partial x^k}=\Gamma_{i,kl}+\Gamma_{l,ik},
  \quad
  -\frac{\partial g_{kl}}{\partial x^i}=-\Gamma_{l,ki}-\Gamma_{k,li}.
\end{equation*}
取这些方程之和的一半，注意到 $\Gamma_{i,kl}=\Gamma_{i,lk}$，我们就可求得
\begin{equation}
  \Gamma_{i,kl}=\frac{1}{2}\left(\frac{\partial g_{ik}}{\partial x^l}
  +\frac{\partial g_{il}}{\partial x^k}
  -\frac{\partial g_{kl}}{\partial x^i}\right).
\end{equation}
由此我们得到符号 $\Gamma^i_{kl}=g^{im}\Gamma_{m,kl}$ 的表达式如下：
\begin{equation}
  \Gamma^i_{kl}=\frac{1}{2}g^{im}\left(\frac{\partial g_{mk}}{\partial x^l}
  +\frac{\partial g_{ml}}{\partial x^k}
  -\frac{\partial g_{kl}}{\partial x^m}\right).
\end{equation}
这些公式就是我们所要求的用度规张量表示克里斯托夫符号的式子.

现在我们来推导对今后有用的缩并克里斯托夫符号 $\Gamma^i_{ki}$ 的表达式.为此，我们计算由张量 $g_{ik}$ 的分量所构成的行列式 $g$ 的微分 $\mathrm{d}g$；取张量 $g_{ik}$ 的每个分量的微分，乘以其在行列式中的系数，亦即乘以相应的子行列式，就可得到 $\mathrm{d}g$.另一方面，大家知道，与 $g_{ik}$ 互逆的张量 $g^{ik}$ 的分量，等于 $g_{ik}$ 的子行列式除以其行列式.因此行列式 $g$ 的子行列式等于 $g\,g^{ik}$.由此可见，
\begin{equation}
  \mathrm{d}g=gg^{ik}\mathrm{d}g_{ik}
  =-gg_{ik}\mathrm{d}g^{ik}
\end{equation}
（既然 $g_{ik}g^{ik}=\delta^i_i=4$，那么，$g^{ik}\mathrm{d}g_{ik}=-g_{ik}\mathrm{d}g^{ik}$）.

从（86.3）式，得到
\begin{equation*}
  \Gamma^i_{ki}=\frac{1}{2}g^{im}\left(\frac{\partial g_{mk}}{\partial x^i}
  +\frac{\partial g_{mi}}{\partial x^k}
  -\frac{\partial g_{ki}}{\partial x^m}\right).
\end{equation*}
将括号中第一项和第三项的指标 m 和 i 改变位置，我们看出，这两项互相抵消，所以
\begin{equation*}
  \Gamma^i_{ki}=\frac{1}{2}g^{im}\frac{\partial g_{im}}{\partial x^k},
\end{equation*}
或者，按照（86.4）
\begin{equation}
  \Gamma^i_{ki}=\frac{1}{2g}\frac{\partial g}{\partial x^k}
  =\frac{\partial\ln\sqrt{-g}}{\partial x^k}.
\end{equation}

写出 $g^{kl}\Gamma^i_{kl}$ 这个量的表达式是有益的；我们有
\begin{equation*}
  g^{kl}\Gamma^i_{kl}=\frac{1}{2}g^{kl}g^{im}\left(
  \frac{\partial g_{mk}}{\partial x^l}
  +\frac{\partial g_{lm}}{\partial x^k}
  -\frac{\partial g_{kl}}{\partial x^m}\right)
  =g^{kl}g^{im}\left(\frac{\partial g_{mk}}{\partial x^l}
  -\frac{1}{2}\frac{\partial g_{kl}}{\partial x^m}\right).
\end{equation*}
利用(86.4)，这个式子可以化为
\begin{equation}
  g^{kl}\Gamma^i_{kl}=-\frac{1}{\sqrt{-g}}
  \frac{\partial(\sqrt{-g}\,g^{ik})}{\partial x^k}.
\end{equation}

为了以后的各种运算，我们要记住逆变张量 $g^{ik}$ 的导数与 $g_{ik}$ 的导数有以下的关系：
\begin{equation}
  g_{il}\frac{\partial g^{lk}}{\partial x^m}
  =-g^{lk}\frac{\partial g_{il}}{\partial x^m}
\end{equation}
（通过将 $g_{il}g^{lk}=\delta^k_l$ 微分而得到）.最后我们指出，$g^{ik}$ 的导数也可以用 $\Gamma^i_{kl}$ 各量来表示，从恒等式 $g^{ik}{}_{;l}=0$ 直接得出
\begin{equation}
  \frac{\partial g^{ik}}{\partial x^l}
  =-\Gamma^i_{ml}g^{mk}-\Gamma^k_{ml}g^{im}.
\end{equation}

将矢量的散度概念推广到曲线坐标中，并利用已经求出的公式，我们可以用便利的形式写出散度 $A^i{}_{;i}$ 的表达式.利用（86.5）式，我们有，
\begin{equation*}
  A^i{}_{;i}=\frac{\partial A^i}{\partial x^i}
  +\Gamma^i_{li}A^l
  =\frac{\partial A^i}{\partial x^i}
  +A^l\frac{\partial\ln\sqrt{-g}}{\partial x^l},
\end{equation*}
或者，最终得到，
\begin{equation}
  A^i{}_{;i}=\frac{1}{\sqrt{-g}}
  \frac{\partial(\sqrt{-g}\,A^i)}{\partial x^i}.
\end{equation}

对于反对称张量 $A^{ik}$ 的导数，我们可以推出类似的表达式.从（85.12)，我们有
\begin{equation*}
  A^{ik}{}_{;k}=\frac{\partial A^{ik}}{\partial x^k}
  +\Gamma^i_{mk}A^{mk}+\Gamma^k_{mk}A^{im}.
\end{equation*}
但是既然 $A^{mk}=-A^{km}$，所以
\begin{equation*}
  \Gamma^i_{mk}A^{mk}=-\Gamma^i_{km}A^{km}=0.
\end{equation*}
将 $\Gamma^k_{mk}$ 用（86.5）式替换，结果我们得到
\begin{equation}
  A^{ik}{}_{;k}=\frac{1}{\sqrt{-g}}
  \frac{\partial(\sqrt{-g}\,A^{ik})}{\partial x^k}.
\end{equation}

现在假设 $A_{ik}$ 是对称张量；我们来计算它的混合分量 $A^k_{i;k}$ 的表达式.我们有
\begin{equation*}
  A^k_{i;k}=\frac{\partial A^k_i}{\partial x^k}
  +\Gamma^k_{lk}A^l_i-\Gamma^l_{ik}A^k_l
  =\frac{1}{\sqrt{-g}}\frac{\partial(A^k_i\sqrt{-g})}{\partial x^k}
  -\Gamma^l_{ki}A^k_l.
\end{equation*}
其中的最后一项等于
\begin{equation*}
  -\frac{1}{2}\left(\frac{\partial g_{il}}{\partial x^k}
  +\frac{\partial g_{kl}}{\partial x^i}
  -\frac{\partial g_{ik}}{\partial x^l}\right)A^{kl}.
\end{equation*}
由于张量 $A^{kl}$ 的对称性，括号内有两项互相抵消，于是留下
\begin{equation}
  A^k_{i;k}=\frac{1}{\sqrt{-g}}
  \frac{\partial(\sqrt{-g}\,A^k_i)}{\partial x^k}
  -\frac{1}{2}\frac{\partial g_{kl}}{\partial x^i}A^{kl}.
\end{equation}

在笛卡儿坐标中，$\frac{\partial A_i}{\partial x^k}-\frac{\partial A_k}{\partial x^i}$ 是一个反对称张量.在曲线坐标中，这个张量是 $A_{i;k}-A_{k;i}$.但是，借助于 $A_{i;k}$ 的表达式，而且由于 $\Gamma^i_{kl}=\Gamma^i_{lk}$，我们有
\begin{equation}
  A_{i;k}-A_{k;i}=\frac{\partial A_i}{\partial x^k}
  -\frac{\partial A_k}{\partial x^i}.
\end{equation}

最后，我们将标量 $\varphi$ 的二阶导数的和 $\frac{\partial^2\varphi}{\partial x_i\partial x^i}$ 变换到曲线坐标中去.显而易见，在曲线坐标中，这个和化为 $\varphi^{;i}{}_{;i}$，但是 $\varphi_{;i}=\partial\varphi/\partial x^i$，因为一个标量的协变微分就是普通微分.将指标上移，我们有
\begin{equation*}
  \varphi^{;i}=g^{ik}\frac{\partial\varphi}{\partial x^k},
\end{equation*}
利用公式（86.9)，我们求得
\begin{equation}
  \varphi^{;i}{}_{;i}=\frac{1}{\sqrt{-g}}
  \frac{\partial}{\partial x^i}\left(\sqrt{-g}\,g^{ik}
  \frac{\partial\varphi}{\partial x^k}\right).
\end{equation}

值得注意的是，一个矢量沿一超曲面的积分变为在一个四维体积上的积分的高斯定理（83.17)，按照（86.9）可以写为
\begin{equation}
  \oint A^i\sqrt{-g}\,\mathrm{d}S_i
  =\int A^i{}_{;i}\sqrt{-g}\,\mathrm{d}\varOmega.
\end{equation}

\section{引力场中粒子的运动}

在狭义相对论中，一个自由粒子的运动由最小作用量原理
\begin{equation}
  \delta S=-mc\,\delta\int\mathrm{d}s=0
\end{equation}
来决定，按照这个原理，粒子是这样运动的，它的世界线在两个给定世界点之间是一条极值曲线；在我们的情形下，这是一条直线（在普通三维空间，它与匀速直线运动相应）.

显而易见，一个粒子在引力场中的运动为与(87.1)式同样的最小作用量原理所决定，因为引力场不是别的，只是四维空间度规的改变，而这个改变也只是表现在 ds 用 $\mathrm{d}x^i$ 表示的式子内的变化而已.因此，在引力场中，粒子是这样运动的，它的世界点沿着极值曲线运动，或者如通常所说，沿着四维空间 $x^0,x^1,x^2,x^3$ 中的测地线运动；然而，因为在存在引力场的情况下，时空不是伽利略的，因而这条线并不“直”，粒子的真实空间运动就既不匀速，也不是直线的了.

作为再次从最小作用量原理出发（见本节习题）的替代方案，更简便的办法是通过对狭义相对论中粒子的自由运动的微分方程进行相应扩展来找到粒子在引力场中的运动方程.在伽利略四维坐标系中，一个自由粒子的运动方程是 $\mathrm{d}u^i/\mathrm{d}s=0$ 或 $\mathrm{d}u^i=0$，此处的 $u^i=\mathrm{d}x^i/\mathrm{d}s$ 是四维速度.显然，在曲线坐标中，这个方程可以推广为方程
\begin{equation}
  \mathrm{D}u^i=0.
\end{equation}

从矢量的协变微分的表达式（85.6）我们可得
\begin{equation*}
  \mathrm{d}u^i+\Gamma^i_{kl}u^k\mathrm{d}x^l=0.
\end{equation*}
用 ds 除这个方程，我们求得
\begin{equation}
  \frac{\mathrm{d}^2x^i}{\mathrm{d}s^2}
  +\Gamma^i_{kl}\frac{\mathrm{d}x^k}{\mathrm{d}s}
  \frac{\mathrm{d}x^l}{\mathrm{d}s}=0.
\end{equation}
这就是要求的运动方程.我们看出，一个粒子在引力场中的运动取决于 $\Gamma^i_{kl}$ 各量.导数 $\mathrm{d}^2x^i/\mathrm{d}s^2$ 是粒子的四维加速度.因此，我们可以称 $-m\Gamma^i_{kl}u^ku^l$ 这个量为“四维力”，这个力作用在位于引力场内的粒子上.这时，张量 $g_{ik}$ 起引力场的“势”的作用——它的导数决定场的“强度”$\Gamma^i_{kl}$\footnote{同时我们给出通过四维加速度的协变分量表示的运动方程的形式. 从条件 $\mathrm{D}u_i=0$ 得出
\begin{equation*}
  \frac{\mathrm{d}u_i}{\mathrm{d}s}-\Gamma_{k,il}u^ku^l=0.
\end{equation*}
将（86.2）中的 $\Gamma_{k,il}$ 代入上式，两项相消，于是得到
\begin{equation*}\tag{87.3a}
  \frac{\mathrm{d}u_i}{\mathrm{d}s}
  -\frac{1}{2}\frac{\partial g_{kl}}{\partial x^i}u^ku^l=0.
\end{equation*}}.

在§85中曾经展示，通过选择相应的坐标系，总是可以将任意指定的时空点上的所有 $\Gamma^i_{kl}$ 归零.现在我们可以看到，选择这样的局部惯性参考系意味着在指定的无限小的时空单元内消除引力场，而这种可能性就是相对论引力场论中等效原理的表达\footnote{在§85最后一个脚注中也提到选取“沿给定的世界线的惯性”参考系的可能性. 特殊情况下，如果这条线是时间坐标线（沿着它 $x^1,x^2,x^3=\mathrm{const}$），那么在给定的空间体积元内，引力场在任何时候都将被消除.}.

引力场中粒子的四维动量还像以前那样定义：
\begin{equation}
  p^i=mcu^i,
\end{equation}
它的平方是
\begin{equation}
  p_ip^i=m^2c^2.
\end{equation}
用 $-\partial S/\partial x^i$ 代替 $p_i$，我们便得到粒子在引力场中的哈密顿-雅可比方程：
\begin{equation}
  g^{ik}\frac{\partial S}{\partial x^i}
  \frac{\partial S}{\partial x^k}-m^2c^2=0.
\end{equation}

(87.3)形式的测地方程，对于光信号的传播是不能应用的，因为我们知道，沿着光线传播的世界线，间隔 $\mathrm{d}s=0$，所以方程（87.3）内所有的项都变为无穷大.为了得到在这种情况下所需要形式的运动方程，我们利用在几何光学中，光线传播的方向决定于与光线相切的波矢量这个事实.因此，我们可以将四维波矢写成 $k^i=\mathrm{d}x^i/\mathrm{d}\lambda$ 的形式，此处的 λ 为某个沿光线变化的参数.在狭义相对论中，当光在真空中传播时，波矢量沿着光线不改变，亦即 $\mathrm{d}k^i=0$（见§53).在引力场中，这个方程显然化为 $\mathrm{D}k^i=0$ 或
\begin{equation}
  \frac{\mathrm{d}k^i}{\mathrm{d}\lambda}
  +\Gamma^i_{kl}k^kk^l=0
\end{equation}
（参数 λ 也由这个方程决定)\footnote{如果沿着测地线 $\mathrm{d}s\equiv 0$，那么这些测地线称为类光测地线或者各向同性测地线.}.

我们知道（见§48)，四维波矢的平方值是零，亦即
\begin{equation}
  k_ik^i=0.
\end{equation}
在此用 $\partial\psi/\partial x^i$ 代替 $k_i$（$\psi$ 是程函)，我们求得引力场中的程函方程如下：
\begin{equation}
  g^{ik}\frac{\partial\psi}{\partial x^i}
  \frac{\partial\psi}{\partial x^k}=0.
\end{equation}

在低速极限情形下，一个粒子在引力场中的相对论运动方程，应当过渡到相应的非相对论的方程.这时，我们必须记着，速度低的假设就已经要求引力场的本身是弱的，假如不是如此，其中的粒子就会有高的速度.

我们来研究在这个极限情形下决定场的度规张量 $g_{ik}$ 与引力场中的非相对论的势 $\varphi$ 的关系如何.

在非相对论力学中，粒子在引力场中的运动，由拉格朗日量(81.1)所决定.加上一个常数 $-mc^2$\footnote{当然，对势 $\varphi$ 的定义可以相差一任意常数. 我们在任何地方都以自然的方式选取这个常数，即使得远离场源物体之处的势变为零.}，我们现在将它写成
\begin{equation}
  L=-mc^2+\frac{mv^2}{2}-m\varphi,
\end{equation}
因此，一个粒子在引力场中的非相对论作用量 S 有下面的形式：
\begin{equation*}
  S=\int L\,\mathrm{d}t
  =-mc\int\left(c-\frac{v^2}{2c}+\frac{\varphi}{c}\right)\mathrm{d}t.
\end{equation*}
将它与相对论的作用量 $S=-mc\int\mathrm{d}s$ 相比较，我们看出，在所考虑的极限情形下，
\begin{equation*}
  \mathrm{d}s=\left(c-\frac{v^2}{2c}+\frac{\varphi}{c}\right)\mathrm{d}t.
\end{equation*}
取平方，略去 $c\to\infty$ 时趋于 0 的项，我们求得
\begin{equation}
  \mathrm{d}s^2=(c^2+2\varphi)\mathrm{d}t^2-\mathrm{d}\boldsymbol{r}^2,
\end{equation}
其中已经用了关系 $\boldsymbol{v}\,\mathrm{d}t=\mathrm{d}\boldsymbol{r}$.

因此，在极限情形下，度规张量的分量 $g_{00}$ 等于
\begin{equation}
  g_{00}=1+\frac{2\varphi}{c^2}.
\end{equation}

关于其他分量，从（87.11）推出 $g_{\alpha\beta}=\delta_{\alpha\beta}$，$g_{0\alpha}=0$.然而，实际上，一般来说，对它们的修正与对 $g_{00}$ 的修正同数量级（关于这一点的详情，请参看§106).之所以不能用上面的方法来决定这些修正是与这个事实有关的，即对 $g_{\alpha\beta}$ 的修正虽然与对 $g_{00}$ 的修正同数量级，然而这个修正在拉格朗日量内产生更高级小量的一些项（因为在 $\mathrm{d}s^2$ 的式子中，分量 $g_{\alpha\beta}$ 不乘以 $c^2$，而 $g_{00}$ 却要乘以 $c^2$）.

\begin{xiti}

\noindent{\heiti 习题}\quad 从最小作用量原理 (87.1) 出发推导运动方程 (87.3).

解：我们有
\begin{equation*}
  \delta\mathrm{d}s^2=2\mathrm{d}s\,\delta\mathrm{d}s
  =\delta(g_{ik}\mathrm{d}x^i\mathrm{d}x^k)
  =\mathrm{d}x^i\mathrm{d}x^k\frac{\partial g_{ik}}{\partial x^l}\delta x^l
  +2g_{ik}\mathrm{d}x^i\mathrm{d}\delta x^k.
\end{equation*}
因此
\begin{equation*}
\begin{aligned}
  \delta S={}&-mc\int\left\{\frac{1}{2}\frac{\mathrm{d}x^i}{\mathrm{d}s}
  \frac{\mathrm{d}x^k}{\mathrm{d}s}\frac{\partial g_{ik}}{\partial x^l}\delta x^l
  +g_{ik}\frac{\mathrm{d}x^i}{\mathrm{d}s}
  \frac{\mathrm{d}\delta x^k}{\mathrm{d}s}\right\}\mathrm{d}s\\
  ={}&-mc\int\left\{\frac{1}{2}\frac{\mathrm{d}x^i}{\mathrm{d}s}
  \frac{\mathrm{d}x^k}{\mathrm{d}s}\frac{\partial g_{ik}}{\partial x^l}\delta x^l
  -\frac{\mathrm{d}}{\mathrm{d}s}\left(g_{ik}\frac{\mathrm{d}x^i}{\mathrm{d}s}\right)
  \delta x^k\right\}\mathrm{d}s
\end{aligned}
\end{equation*}
（在进行分部积分的时候考虑到积分范围两端 $\delta x^k=0$）.在第二个被积分项中用指标 l 替换指标 k.那么，在对 $\delta x^l$ 进行任意变分时令系数等于零，我们得到
\begin{equation*}
  \frac{1}{2}u^iu^k\frac{\partial g_{ik}}{\partial x^l}
  -\frac{\mathrm{d}}{\mathrm{d}s}(g_{il}u^i)
  =\frac{1}{2}u^iu^k\frac{\partial g_{ik}}{\partial x^l}
  -g_{il}\frac{\mathrm{d}u^i}{\mathrm{d}s}
  -u^iu^k\frac{\partial g_{il}}{\partial x^k}=0.
\end{equation*}
这里指出，第三项可以写成如下形式
\begin{equation*}
  -\frac{1}{2}u^iu^k\left(\frac{\partial g_{il}}{\partial x^k}
  +\frac{\partial g_{kl}}{\partial x^i}\right),
\end{equation*}
引入克里斯托夫符号 $\Gamma_{l,ik}$，按照（86.2）我们得到
\begin{equation*}
  g_{il}\frac{\mathrm{d}u^i}{\mathrm{d}s}
  +\Gamma_{l,ik}u^iu^k=0.
\end{equation*}
通过上移指标 l 即可得到方程 (87.3).

\end{xiti}

\section{恒定引力场}

如果能够选择参考系，使其中度规张量的所有分量均不依赖于时间坐标 $x^0$，能满足这样条件的引力场称为恒定引力场；时间坐标 $x^0$ 称做世界时间.

世界时间的选择不是唯一的，如果给 $x^0$ 加上空间坐标的一个任意函数，所有的 $g_{ik}$ 仍旧没有包含 $x^0$；这个变换相应于在空间每一点选择时间原点的任意性\footnote{容易看到，在这个变换下，空间度规应该是不变的. 实际上，利用任意的函数 $f(x^1,x^2,x^3)$ 作如下替换
\begin{equation*}
  x^0\to x^0+f(x^1,x^2,x^3),
\end{equation*}
分量 $g_{ik}$ 按照下列各式替换
\begin{equation*}
\begin{aligned}
  &g_{\alpha\beta}\to g_{\alpha\beta}+g_{00}f_{,\alpha}f_{,\beta}
  -g_{0\alpha}f_{,\beta}-g_{0\beta}f_{,\alpha},\\
  &g_{0\alpha}\to g_{0\alpha}-g_{00}f_{,\alpha},\qquad g_{00}\to g_{00},
\end{aligned}
\end{equation*}
其中 $f_{,\alpha}\equiv\partial f/\partial x^{\alpha}$. 在此条件下，显然，三维张量（84.7）保持不变.}.此外，世界时间还可以乘以一个任意常数，即世界时间测量单位的选择是任意的.

严格地说，只有一个物体产生的引力场才能是恒定的.在多个物体构成的系统中，它们互相之间的吸引会导致运动，其结果是它们产生的引力场不可能恒定.

如果产生引力场的物体是静止的（在一个 $g_{ik}$ 与 $x^0$ 无关的参考系内)，这时，两个时间方向是等价的.在空间所有点适当选择时间原点，间隔 ds 在改变 $x^0$ 的符号时应当不变.从此可以断定，在这种情形下，度规张量 $g_{0\alpha}$ 的所有分量都为零.这类恒定引力场称为静态引力场.

但是物体的静止状态并不是产生恒定引力场的必要条件.围绕自己的对称轴作匀速旋转的轴对称物体产生的引力场同样是恒定的.但在这种情形下，两个时间方向绝不是等价的；假如时间的符号改变，那么，例如，旋转的角速度的符号也将会改变.显而易见，在这种类型的恒定引力场中，度规张量的分量 $g_{0\alpha}$ 一般来说并不为零.我们称这一类恒定场为稳态引力场.

在恒定引力场中世界时间的意义在于：在空间某一点发生的两个事件的世界时间间隔，同发生在空间另一点的任何另外两个事件之间的世界时间间隔相等，如果这些事件分别与第一对事件（在§84解释的意义上）是同时的.但相同的世界时间 $x^0$ 的间隔在空间的不同点上对应着不同的固有时 $\tau$ 的间隔.世界时间与固有时的关系，即公式（84.1）现在可以写成
\begin{equation}
  \tau=\frac{1}{c}\sqrt{g_{00}}\,x^0,
\end{equation}
这个关系式可以应用于任何有限时间间隔.

假如引力场是弱的，那么，我们可以用近似式（87.12）和（88.1）得到
\begin{equation}
  \tau=\frac{x^0}{c}\left(1+\frac{\varphi}{c^2}\right).
\end{equation}
因此，固有时流逝得愈慢，空间一给定点的引力势就愈小，也就是说引力势的绝对值愈大（以后在§96中，可以证明 $\varphi$ 是负的).假如将两个完全一样的钟之一置于引力场内，那么，在引力场内的钟要走得慢些.

上面已经指出，在静态引力场中，度规张量的分量 $g_{0\alpha}$ 是零.按照§84 的结果，这就意味着，在这样一个场中，钟的同步在整个空间都是可能的.我们也应注意，在静态场中，空间距离元不过是
\begin{equation}
  \mathrm{d}l^2=-g_{\alpha\beta}\mathrm{d}x^{\alpha}\mathrm{d}x^{\beta}.
\end{equation}

在稳态场中，$g_{0\alpha}$ 不等于零，在整个空间中时钟的同步是不可能的.既然 $g_{ik}$ 与 $x^0$ 无关，那么在空间不同点发生的两个同时事件的世界时间之差的公式（84.14）可以写成如下形式：
\begin{equation}
  \Delta x^0=-\int\frac{g_{0\alpha}\mathrm{d}x^{\alpha}}{g_{00}},
\end{equation}
当沿着一条线同步钟时，式(88.4)对于这条线上的任意两点都适用.当沿着一条闭合回路同步钟时，世界时间之差(它在回到出发点时应该被记录下来)等于沿着闭合回路所取的积分\footnote{如果和式 $g_{0\alpha}\mathrm{d}x^{\alpha}/g_{00}$ 是空间坐标的某一函数的全微分，则积分（88.5）恒等于零. 然而这种情形只意味着，我们实际上处理的是静态场，通过形如 $x^0\to x^0+f(x^{\alpha})$ 的变换总是可以使所有 $g_{0\alpha}$ 变为零.}
\begin{equation}
  \Delta x^0=-\oint\frac{g_{0\alpha}\mathrm{d}x^{\alpha}}{g_{00}}.
\end{equation}

让我们来考虑一条光线在一恒定引力场内的传播.在§53中，我们已经知道，光的频率是程函 $\psi$ 的时间导数（符号相反).因此，用世界时间 $x^0/c$ 表示的频率是 $\omega_0=-c\,\partial\psi/\partial x^0$.既然恒定场中的程函方程（87.9)不显含 $x^0$，那么，在光线传播时，频率 $\omega_0$ 保持不变.用固有时来测量的频率是 $\omega=-\partial\psi/\partial\tau$；这个频率在空间的不同点是不同的.

由于关系式
\begin{equation*}
  \frac{\partial\psi}{\partial\tau}
  =\frac{\partial\psi}{\partial x^0}\frac{\partial x^0}{\partial\tau}
  =\frac{\partial\psi}{\partial x^0}\frac{c}{\sqrt{g_{00}}},
\end{equation*}
我们就有
\begin{equation}
  \omega=\frac{\omega_0}{\sqrt{g_{00}}}.
\end{equation}
在弱引力场中，由此可近似地得到
\begin{equation}
  \omega=\omega_0\left(1-\frac{\varphi}{c^2}\right).
\end{equation}

我们看出，光的频率随着引力场的势的绝对值的增加而增大，也就是说，当光射向产生场的物体时，频率就随之而增加；反之，当光离开物体时，频率随之而减小.假如一条光线从一点射出，而这一点的引力势为 $\varphi_1$，光线（在这一点）的频率为 $\omega$，那么，当到达引力势为 $\varphi_2$ 的另外一点时，光线的频率(以在那一点的固有时来测量)就等于
\begin{equation*}
  \frac{\omega}{1-\dfrac{\varphi_1}{c^2}}
  \left(1-\frac{\varphi_2}{c^2}\right)
  \approx\omega\left(1+\frac{\varphi_1-\varphi_2}{c^2}\right).
\end{equation*}
例如，在太阳上观察到的由太阳上的原子所发出的光谱线，与在地球上观察到的由地球上的同样原子所发出的光谱线是一样的，假如我们在地球上观察太阳上的原子所发出的光谱，那么，根据上面所说的，光谱线对于地球上的同样原子所发射出的光谱线而言就会有些移动.每条以 $\omega$ 为频率的光线将移动一间隔 $\Delta\omega$，而 $\Delta\omega$ 由公式
\begin{equation}
  \Delta\omega=\frac{\varphi_1-\varphi_2}{c^2}\,\omega
\end{equation}
所决定，其中的 $\varphi_1$ 和 $\varphi_2$ 分别是光谱发射点和观察点的引力场的势.假如我们在地球上观察由太阳或恒星所发出的光谱，那么，$\left|\varphi_1\right|>\left|\varphi_2\right|$，从（88.8）可以断定 $\Delta\omega<0$，就是说，移向频率减小的方向.上述的现象称为红移.

根据前面关于世界时间的讲述，可以直接解释发生这个现象的原因.因为场是恒定的，光波从空间中一给定点传播到另一点的过程中，某振动所需的世界时间间隔与时间 $x^0$ 无关.因此很明显，在单位世界时间间隔内所发生的振动数在光线上所有的点都是一样的.但是，对于同一个世界时间间隔，我们离开产生场的物体愈远，与之相对应的固有时间隔就愈大.所以，当光从这些物质离开时，频率将要降低，就是说，每单位固有时内的振动数将要减少.

当粒子在恒定场中运动时，它的能量是守恒的，这个能量是用作用量对于世界时间的导数 $\left(-c\,\dfrac{\partial S}{\partial x^0}\right)$ 来定义的；从 $x^0$ 并不明显地出现在哈密顿-雅可比方程中这一事实就可推出该结论.这样定义的能量是协变四维动量矢量 $p_k=mcu_k=mcg_{ki}u^i$ 的时间分量.在静态场中，$\mathrm{d}s^2=g_{00}(\mathrm{d}x^0)^2-\mathrm{d}l^2$，因而对于能量（我们用 $\mathcal{E}_0$ 表示），我们有
\begin{equation*}
  \mathcal{E}_0=mc^2g_{00}\frac{\mathrm{d}x^0}{\mathrm{d}s}
  =mc^2g_{00}\frac{\mathrm{d}x^0}
  {\sqrt{g_{00}(\mathrm{d}x^0)^2-\mathrm{d}l^2}}.
\end{equation*}
我们引入用固有时来测量的，亦即由在给定点的观察者来测量的粒子速度，
\begin{equation*}
  v=\frac{\mathrm{d}l}{\mathrm{d}\tau}
  =\frac{c\,\mathrm{d}l}{\sqrt{g_{00}}\,\mathrm{d}x^0},
\end{equation*}
则我们得到能量
\begin{equation}
  \mathcal{E}_0=\frac{mc^2\sqrt{g_{00}}}
  {\sqrt{1-\dfrac{v^2}{c^2}}}.
\end{equation}
这是在粒子运动时守恒的量.

容易证明，能量的表达式(88.9)在恒定场中也成立，只要速度 v 以固有时测量，固有时由沿粒子轨道同步的时钟确定.如果粒子在世界时间的 $x^0$ 时刻从 A 点出发，并且在 $x^0+\mathrm{d}x^0$ 时刻到达无限接近的 B 点，那么现在为了确定速度就不应该取时间段 $(x^0+\mathrm{d}x^0)-x^0=\mathrm{d}x^0$，而应该取 $x^0+\mathrm{d}x^0$ 和时刻 $x^0-\dfrac{g_{0\alpha}}{g_{00}}\mathrm{d}x^{\alpha}$ 的差值，后者在 B 点与 A 点的 $x^0$ 时刻是同时的：
\begin{equation*}
  (x^0+\mathrm{d}x^0)
  -\left(x^0-\frac{g_{0\alpha}}{g_{00}}\mathrm{d}x^{\alpha}\right)
  =\mathrm{d}x^0+\frac{g_{0\alpha}}{g_{00}}\mathrm{d}x^{\alpha}.
\end{equation*}
给上式乘以 $\sqrt{g_{00}}/c$，我们得到相应的固有时间隔，则速度是
\begin{equation}
  v^{\alpha}=\frac{c\,\mathrm{d}x^{\alpha}}
  {\sqrt{h}\,(\mathrm{d}x^0-g_{\alpha}\mathrm{d}x^{\alpha})},
\end{equation}
其中我们引入了表示符号：
\begin{equation}
  g_{\alpha}=-\frac{g_{0\alpha}}{g_{00}},\qquad h=g_{00}
\end{equation}
来描述三维矢量 $\boldsymbol{g}$（已经在§84 中提到过）和三维标量 $g_{00}$.速度 $\boldsymbol{v}$ 的协变分量就像具有度规 $\gamma_{\alpha\beta}$ 的空间中三维矢量的协变分量，而这个矢量的平方应该相应地是\footnote{接下来我们将不止一次地在讨论中与四维矢量和四维张量一起引入三维矢量和张量，后者定义在具有度规 $\gamma_{\alpha\beta}$ 的空间中；已经引入的矢量 $\boldsymbol{g}$ 和 $\boldsymbol{v}$ 就是这种类型的特例. 四维张量运算（包括指标的上移和下移）利用度规张量 $g_{ik}$ 进行，而三维张量运算则利用张量 $\gamma_{\alpha\beta}$. 为了避免由此可能产生的误解，凡是用于表示三维量的符号，我们将不用其表示四维量.}：
\begin{equation}
  v_{\alpha}=\gamma_{\alpha\beta}v^{\beta},\qquad
  v^2=v_{\alpha}v^{\alpha}.
\end{equation}

我们指出，在这样的定义下，间隔 ds 可以通过速度用平常的形式表示出来：
\begin{equation}
\begin{aligned}[b]
  \mathrm{d}s^2&=g_{00}(\mathrm{d}x^0)^2
  +2g_{0\alpha}\mathrm{d}x^0\mathrm{d}x^{\alpha}
  +g_{\alpha\beta}\mathrm{d}x^{\alpha}\mathrm{d}x^{\beta}\\
  &=h(\mathrm{d}x^0-g_{\alpha}\mathrm{d}x^{\alpha})^2-\mathrm{d}l^2
  =h(\mathrm{d}x^0-g_{\alpha}\mathrm{d}x^{\alpha})^2
  \left(1-\frac{v^2}{c^2}\right).
\end{aligned}
\end{equation}
四维速度的分量 $u^i=\mathrm{d}x^i/\mathrm{d}s$ 等于
\begin{equation}
  u^{\alpha}=\frac{v^{\alpha}}{c\sqrt{1-\dfrac{v^2}{c^2}}},\qquad
  u^0=\frac{1}{\sqrt{h}\sqrt{1-\dfrac{v^2}{c^2}}}
  +\frac{g_{\alpha}v^{\alpha}}{c\sqrt{1-\dfrac{v^2}{c^2}}}.
\end{equation}
而能量
\begin{equation*}
  \mathcal{E}_0=mc^2g_{0i}u^i=mc^2h(u^0-g_{\alpha}u^{\alpha})
\end{equation*}
代入（88.14）后得到 (88.9)式.

在弱引力场和低速运动的极限条件下，将 $g_{00}=1+2\varphi/c^2$ 代入（88.9），近似得到
\begin{equation}
  \mathcal{E}_0=mc^2+\frac{mv^2}{2}+m\varphi,
\end{equation}
其中 $m\varphi$ 为粒子在引力场中的势能，与拉格朗日量（87.10）对应.

\begin{xiti}

\noindent{\heiti 习题1}\quad 求恒定引力场中作用在粒子上的力.

解：针对我们需要的分量 $\Gamma^i_{kl}$ 找到下列表达式
\begin{equation*}
\begin{aligned}
  \Gamma^{\alpha}_{00}&=\frac{1}{2}h^{;\alpha},\\
  \Gamma^{\alpha}_{0\beta}
  &=\frac{h}{2}\left(g^{\alpha}{}_{;\beta}
  -g_{\beta}{}^{;\alpha}\right)
  -\frac{1}{2}g_{\beta}h^{;\alpha},\\
  \Gamma^{\alpha}_{\beta\gamma}
  &=\lambda^{\alpha}_{\beta\gamma}
  +\frac{h}{2}\left[g_{\beta}\left(g_{\gamma}{}^{;\alpha}
  -g^{\alpha}{}_{;\gamma}\right)
  +g_{\gamma}\left(g_{\beta}{}^{;\alpha}
  -g^{\alpha}{}_{;\beta}\right)\right]
  +\frac{1}{2}g_{\beta}g_{\gamma}h^{;\alpha}.
\end{aligned}\tag{1}
\end{equation*}
在这些表达式中所有的张量运算（协变微分，指标的上移和下移）都在带有度规 $\gamma_{\alpha\beta}$ 的三维空间里针对三维矢量 $g^{\alpha}$ 和三维标量 h(88.11)进行；$\lambda^{\alpha}_{\beta\gamma}$ 是三维克里斯托夫符号，由张量 $\gamma_{\alpha\beta}$ 的分量构成，如同 $g_{ik}$ 的分量构成 $\Gamma^i_{kl}$ 那样；在计算时使用公式（84.9)—(84.12).

将 (1) 代入运动方程
\begin{equation*}
  \frac{\mathrm{d}u^{\alpha}}{\mathrm{d}s}
  =-\Gamma^{\alpha}_{00}(u^0)^2
  -2\Gamma^{\alpha}_{0\beta}u^0u^{\beta}
  -\Gamma^{\alpha}_{\beta\gamma}u^{\beta}u^{\gamma},
\end{equation*}
对四维速度的分量应用公式(88.14)，经过简单的变换得到
\begin{equation*}
  \frac{\mathrm{d}}{\mathrm{d}s}
  \frac{v^{\alpha}}{c\sqrt{1-\dfrac{v^2}{c^2}}}
  =-\frac{h^{;\alpha}}{2h\left(1-\dfrac{v^2}{c^2}\right)}
  -\frac{\sqrt{h}\left(g^{\alpha}{}_{;\beta}
  -g_{\beta}{}^{;\alpha}\right)v^{\beta}}
  {c\left(1-\dfrac{v^2}{c^2}\right)}
  -\frac{\lambda^{\alpha}_{\beta\gamma}v^{\beta}v^{\gamma}}
  {c^2\left(1-\dfrac{v^2}{c^2}\right)}.
\tag{2}\end{equation*}
作用在粒子上的力 $\boldsymbol{f}$ 是它的动量 $\boldsymbol{p}$ 对（校准的）固有时的导数，可利用三维协变微分来确定：
\begin{equation*}
  f^{\alpha}=c\sqrt{1-\frac{v^2}{c^2}}\,
  \frac{\mathrm{D}p^{\alpha}}{\mathrm{d}s}
  =c\sqrt{1-\frac{v^2}{c^2}}\,\frac{\mathrm{d}}{\mathrm{d}s}
  \frac{mv^{\alpha}}{\sqrt{1-\dfrac{v^2}{c^2}}}
  +\lambda^{\alpha}_{\beta\gamma}
  \frac{mv^{\beta}v^{\gamma}}{\sqrt{1-\dfrac{v^2}{c^2}}}.
\end{equation*}
因此从 (2) 我们得到（为了便利将指标 $\alpha$ 下移）
\begin{equation*}
  f_{\alpha}=\frac{mc^2}{\sqrt{1-\dfrac{v^2}{c^2}}}
  \left\{-\frac{\partial}{\partial x^{\alpha}}\ln\sqrt{h}
  +\sqrt{h}\left(\frac{\partial g_{\beta}}{\partial x^{\alpha}}
  -\frac{\partial g_{\alpha}}{\partial x^{\beta}}\right)
  \frac{v^{\beta}}{c}\right\},
\end{equation*}
或者采用一般的三维的矢量表示方法\footnote{在三维曲线坐标中单位反对称张量如下定义
\begin{equation*}
  \eta_{\alpha\beta\gamma}=\sqrt{\gamma}\,e_{\alpha\beta\gamma},\qquad
  \eta^{\alpha\beta\gamma}=\frac{1}{\sqrt{\gamma}}e^{\alpha\beta\gamma},
\end{equation*}
其中 $e_{123}=e^{123}=1$，而当两个指标互换位置时会改变符号（对比（83.13），（83.14））. 与此相应，与反对称张量 $c_{\beta\gamma}=a_{\beta}b_{\gamma}-a_{\gamma}b_{\beta}$ 对偶的矢量 $\boldsymbol{c}=\boldsymbol{a}\times\boldsymbol{b}$ 具有分量
\begin{equation*}
  c_{\alpha}=\frac{1}{2}\sqrt{\gamma}\,e_{\alpha\beta\gamma}c^{\beta\gamma}
  =\sqrt{\gamma}\,e_{\alpha\beta\gamma}a^{\beta}b^{\gamma},\qquad
  c^{\alpha}=\frac{1}{2\sqrt{\gamma}}e^{\alpha\beta\gamma}c_{\beta\gamma}
  =\frac{1}{\sqrt{\gamma}}e^{\alpha\beta\gamma}a_{\beta}b_{\gamma}.
\end{equation*}
反过来
\begin{equation*}
  c_{\alpha\beta}=\sqrt{\gamma}\,e_{\alpha\beta\gamma}c^{\gamma},\qquad
  c^{\alpha\beta}=\frac{1}{\sqrt{\gamma}}e^{\alpha\beta\gamma}c_{\gamma}.
\end{equation*}
特别地，$\operatorname{rot}\boldsymbol{a}$ 在这个意义上应该理解为和张量
$a_{\beta;\alpha}-a_{\alpha;\beta}
=\dfrac{\partial a_{\beta}}{\partial x^{\alpha}}
-\dfrac{\partial a_{\alpha}}{\partial x^{\beta}}$
对偶的矢量，于是它的逆变分量
\begin{equation*}
  (\operatorname{rot}\boldsymbol{a})^{\alpha}
  =\frac{1}{2\sqrt{\gamma}}e^{\alpha\beta\gamma}
  \left(\frac{\partial a_{\gamma}}{\partial x^{\beta}}
  -\frac{\partial a_{\beta}}{\partial x^{\gamma}}\right).
\end{equation*}
与之相关我们还应提到，矢量的三维散度
\begin{equation*}
  \operatorname{div}\boldsymbol{a}
  =\frac{1}{\sqrt{\gamma}}\frac{\partial}{\partial x^{\alpha}}
  (\sqrt{\gamma}\,a^{\alpha})
\end{equation*}
（对比（86.9））.
和正交曲线坐标中三维矢量运算经常使用的公式（例如，见本教程第八卷附录）进行比较时，为了避免误解我们指出，在这些公式中矢量的分量指的是 $\sqrt{g_{11}}A^1(=\sqrt{A_1A^1})$，$\sqrt{g_{22}}A^2$，$\sqrt{g_{33}}A^3$ 各量.}
\begin{equation*}
  \boldsymbol{f}=\frac{mc^2}{\sqrt{1-\dfrac{v^2}{c^2}}}
  \left\{-\operatorname{grad}\ln\sqrt{h}
  +\sqrt{h}\,\frac{\boldsymbol{v}}{c}
  \times(\operatorname{rot}\boldsymbol{g})\right\}.
\tag{3}\end{equation*}
我们指出，如果物体静止，那么作用在它上面的力（(3)式中的第一项)具有势.在运动速度很小的时候(3)式中的第二项具有类似于科里奥利力的形式 $mc\sqrt{h}\,\boldsymbol{v}\times(\operatorname{rot}\boldsymbol{g})$，后者产生于以如下角速度旋转的坐标系（没有场）中：
\begin{equation*}
  \varOmega=\frac{c}{2}\sqrt{h}\,\operatorname{rot}\boldsymbol{g}.
\end{equation*}

\noindent{\heiti 习题2}\quad 导出光线在恒定引力场中传播的费马原理.

解：费马原理（见§53）的内容是
\begin{equation*}
  \delta\int k_{\alpha}\mathrm{d}x^{\alpha}=0,
\end{equation*}
其中的积分是沿着光线而取的，而被积函数必须以频率 $\omega_0$（$\omega_0$ 沿光线是常数)和坐标的微分来表示.注意到 $k_0=-\partial\psi/\partial x^0=\omega_0/c$，我们可写出
\begin{equation*}
  \frac{\omega_0}{c}=k_0=g_{0i}k^i=g_{00}k^0+g_{0\alpha}k^{\alpha}
  =h(k^0-g_{\alpha}k^{\alpha}).
\end{equation*}
将此式代入 $k_ik^i=g_{ik}k^ik^k=0$，并写成
\begin{equation*}
  h(k^0-g_{\alpha}k^{\alpha})^2-\gamma_{\alpha\beta}k^{\alpha}k^{\beta}=0,
\end{equation*}
我们便得到
\begin{equation*}
  \frac{1}{h}\left(\frac{\omega_0}{c}\right)^2
  -\gamma_{\alpha\beta}k^{\alpha}k^{\beta}=0.
\end{equation*}
又根据矢量 $k^{\alpha}$ 应当与矢量 $\mathrm{d}x^{\alpha}$ 同方向的事实，我们便求得
\begin{equation*}
  k^{\alpha}=\frac{\omega_0}{c\sqrt{h}}
  \frac{\mathrm{d}x^{\alpha}}{\mathrm{d}l},
\end{equation*}
其中，dl(84.6)是沿着光线的空间距离元.为了求 $k_{\alpha}$ 的表达式，我们写出
\begin{equation*}
  k^{\alpha}=g^{\alpha i}k_i=g^{\alpha 0}k_0+g^{\alpha\beta}k_{\beta}
  =-g^{\alpha}\frac{\omega_0}{c}-\gamma^{\alpha\beta}k_{\beta},
\end{equation*}
由此
\begin{equation*}
  k_{\alpha}=-\gamma_{\alpha\beta}\left(k^{\beta}
  +\frac{\omega_0}{c}g^{\beta}\right)
  =-\frac{\omega_0}{c}\left(\frac{\gamma_{\alpha\beta}}{\sqrt{h}}
  \frac{\mathrm{d}x^{\beta}}{\mathrm{d}l}+g_{\alpha}\right).
\end{equation*}
最后，乘之以 $\mathrm{d}x^{\alpha}$，我们得到下面形式的费马原理（略去常数因子 $\omega_0/c$）：
\begin{equation*}
  \delta\int\left(\frac{\mathrm{d}l}{\sqrt{h}}
  +g_{\alpha}\mathrm{d}x^{\alpha}\right)=0.
\end{equation*}
在静态场中，我们简单地得到
\begin{equation*}
  \delta\int\frac{\mathrm{d}l}{\sqrt{h}}=0.
\end{equation*}
请注意这个事实，在引力场中，光线并非沿着空间最短的线传播，因为沿着空间的最短的线应该由方程 $\delta\int\mathrm{d}l=0$ 确定.

\end{xiti}

\section{旋转}

作为稳态引力场的一个特例，我们来考虑匀速旋转参考系.

为了求间隔 ds，我们来做一个从静止（惯性）系到匀速旋转系的变换.在静止的坐标系 $r',\varphi',z',t$（我们使用柱坐标 $r',\varphi',z'$)中，间隔有如下形式：
\begin{equation}
  \mathrm{d}s^2=c^2\mathrm{d}t^2-\mathrm{d}r'^2
  -r'^2\mathrm{d}\varphi'^2-\mathrm{d}z'^2.
\end{equation}

设在旋转系中的柱坐标为 $r,\varphi,z$.假如旋转轴与 z 轴和 $z'$ 轴重合，那么，我们有 $r'=r$，$z'=z$，$\varphi'=\varphi+\varOmega t$，此处的 Ω 为旋转的角速度.将这些式子代入(89.1)，我们得到所求的在旋转参考系中间隔的表达式：
\begin{equation}
  \mathrm{d}s^2=(c^2-\varOmega^2r^2)\mathrm{d}t^2
  -2\varOmega r^2\mathrm{d}\varphi\,\mathrm{d}t
  -\mathrm{d}z^2-r^2\mathrm{d}\varphi^2-\mathrm{d}r^2.
\end{equation}

必须注意，旋转参考系仅仅可以应用到距离转动中心 $c/\varOmega$ 的范围之内.事实上，从（89.2)可以看到，当 $r>c/\varOmega$ 时，$g_{00}$ 变为负值，而这是不允许的.旋转系对于大的距离之所以不能应用是因为在大的距离处，速度将大于光速，因此，这样的参考系不可能用真实的物体来实现.

同在一切稳态场中一样，在旋转物体上的钟不可能在所有点上都被单值地校准.当沿着任何封闭曲线进行钟的校准并回到出发点时，我们得到一个时间，这个时间与原来的时间的差值是（见（88.5)）
\begin{equation*}
  \Delta t=-\frac{1}{c}\oint\frac{g_{0\alpha}}{g_{00}}\mathrm{d}x^{\alpha}
  =\frac{1}{c^2}\oint
  \frac{\varOmega r^2\mathrm{d}\varphi}
  {1-\dfrac{\varOmega^2r^2}{c^2}},
\end{equation*}
假设 $\varOmega r/c\ll1$（即旋转速度比光速小得多)，则差值是
\begin{equation}
  \Delta t=\frac{\varOmega}{c^2}\int r^2\mathrm{d}\varphi
  =\pm\frac{2\varOmega}{c^2}S,
\end{equation}
其中，S 是回路包围的面在垂直于旋转轴的一个平面上的投影面积（符号用 $+$ 或 $-$，依我们顺旋转方向或逆旋转方向行走而定）.

假定有一条光线沿着某一条闭合回路传播.让我们来计算光线从出发到回到原点所经过的时间 t（准确到与 $v/c$ 同数量级的项）.假如沿着这一条闭合曲线时间是校准了的，又假设在每一点上我们都用固有时，那么，根据定义，光的速度将永等于 c.既然固有时与世界时间之差与 $v^2/c^2$ 同数量级，那么，在计算所要求的时间间隔 t 时，若只要求准确到与 $v/c$ 同量级的量，这个差就可以略去不计了.因此，我们有
\begin{equation*}
  t=\frac{L}{c}\pm\frac{2\varOmega}{c^2}S,
\end{equation*}
其中，L 是回路的长度.与此相应，用比值 $L/t$ 来测量的光速等于
\begin{equation}
  c\pm 2\varOmega\frac{S}{L}.
\end{equation}

这个公式，如多普勒效应的一级近似公式一样，也可以很容易地用纯经典的方法求得.

\begin{xiti}

\noindent{\heiti 习题}\quad 求旋转坐标系统中的空间距离元.

解：利用（84.6)，（84.7）我们求得
\begin{equation*}
  \mathrm{d}l^2=\mathrm{d}r^2+\mathrm{d}z^2
  +\frac{r^2\mathrm{d}\varphi^2}{1-\dfrac{\varOmega^2r^2}{c^2}},
\end{equation*}
此式决定旋转参考系中的空间几何.我们注意，在平面 $z=\mathrm{const}$ 内的圆（圆心在转轴上)的周长与其半径 r 之比等于
\begin{equation*}
  \frac{2\pi}{\sqrt{1-\dfrac{\varOmega^2r^2}{c^2}}}>2\pi.
\end{equation*}

\end{xiti}

\section{引力场存在时的电动力学方程}

狭义相对论中的电磁场方程很容易推广，使其在一个任意的四维曲线坐标系中也能应用，就是说，在引力场存在时，也能应用.

在狭义相对论中，电磁场张量的定义是：$F_{ik}=\dfrac{\partial A_k}{\partial x^i}-\dfrac{\partial A_i}{\partial x^k}$.显而易见，电磁场张量现在必须相应地定义为 $F_{ik}=A_{k;i}-A_{i;k}$.但是由于（86.12），
\begin{equation}
  F_{ik}=A_{k;i}-A_{i;k}=\frac{\partial A_k}{\partial x^i}
  -\frac{\partial A_i}{\partial x^k},
\end{equation}
因此，$F_{ik}$ 和势 $A_k$ 的关系不改变.由于这个原因，第一对麦克斯韦方程(26.5)也不改变它们的形式\footnote{容易发现，这个方程同时还可以写为形式：
\begin{equation*}
  F_{ik;l}+F_{li;k}+F_{kl;i}=0,
\end{equation*}
由此它的协变性是显而易见的.}：
\begin{equation}
  \frac{\partial F_{ik}}{\partial x^l}
  +\frac{\partial F_{li}}{\partial x^k}
  +\frac{\partial F_{kl}}{\partial x^i}=0.
\end{equation}

为了写出第二对麦克斯韦方程，首先我们必须决定在曲线坐标中的四维电流矢量.这可以用完全与§28相同的办法完成.由空间坐标元 $\mathrm{d}x^1,\mathrm{d}x^2,\mathrm{d}x^3$ 构成的体积元为 $\sqrt{\gamma}\,\mathrm{d}V$，其中 $\gamma$ 是空间度规张量（84.7)的行列式，而 $\mathrm{d}V=\mathrm{d}x^1\mathrm{d}x^2\mathrm{d}x^3$（见§83 最后一个脚注).通过定义 $\mathrm{d}e=\rho\sqrt{\gamma}\,\mathrm{d}V$ 引入电荷密度 $\rho$，其中 de 是体积元 $\sqrt{\gamma}\,\mathrm{d}V$ 内的电荷.给这个等式两边同乘以 $\mathrm{d}x^i$，我们有：
\begin{equation*}
  \mathrm{d}e\,\mathrm{d}x^i=\rho\,\mathrm{d}x^i\sqrt{\gamma}\,
  \mathrm{d}x^1\mathrm{d}x^2\mathrm{d}x^3
  =\frac{\rho}{\sqrt{g_{00}}}\sqrt{-g}\,\mathrm{d}\varOmega
  \frac{\mathrm{d}x^i}{\mathrm{d}x^0}
\end{equation*}
（这里我们使用了公式 $-g=\gamma g_{00}$（84.10)）.乘积 $\sqrt{-g}\,\mathrm{d}\varOmega$ 是不变的四维体积元，因此四维电流矢量等于
\begin{equation}
  j^i=\frac{\rho c}{\sqrt{g_{00}}}\frac{\mathrm{d}x^i}{\mathrm{d}x^0}
\end{equation}
（其中，$\mathrm{d}x^i/\mathrm{d}x^0$ 是坐标随“时间”$x^0$ 变化的速度，其本身并不是一个四维矢量!).四维电流矢量的分量 $j^0$ 乘以 $\sqrt{g_{00}}/c$ 是电荷的空间密度.

对于点电荷而言，电荷密度 $\rho$ 表示为 δ 函数的和，类似于公式(28.1).然而，在这种情况下，应该修正这些函数在曲线坐标条件下的定义.我们将仍按照以前那样把 $\delta(\boldsymbol{r})$ 理解为乘积 $\delta(x^1)\delta(x^2)\delta(x^3)$，而不考虑坐标 $x^1,x^2,x^3$ 的几何意义；那么在 $\mathrm{d}V$（而不是 $\sqrt{\gamma}\,\mathrm{d}V$）上的积分等于1: $\int\delta(\boldsymbol{r})\mathrm{d}V=1$.按照 δ 函数同样的定义，电荷密度是
\begin{equation*}
  \rho=\sum_a\frac{e_a}{\sqrt{\gamma}}\delta(\boldsymbol{r}-\boldsymbol{r}_a),
\end{equation*}
而四维电流矢量是
\begin{equation}
  j^i=\sum_a\frac{e_ac}{\sqrt{-g}}\delta(\boldsymbol{r}-\boldsymbol{r}_a)
  \frac{\mathrm{d}x^i}{\mathrm{d}x^0}.
\end{equation}

电荷守恒由连续性方程来表达，与(29.4)的不同之处仅在于将通常的微分替换为协变微分：
\begin{equation}
  j^i{}_{;i}=\frac{1}{\sqrt{-g}}
  \frac{\partial}{\partial x^i}\left(\sqrt{-g}\,j^i\right)=0
\end{equation}
（使用了公式（86.9)）.

采用类似的方法可以对麦克斯韦方程组(30.2)的第二对进行推广；将其中的通常微分替换为协变微分，我们得到：
\begin{equation}
  F^{ik}{}_{;k}=\frac{1}{\sqrt{-g}}
  \frac{\partial}{\partial x^k}\left(\sqrt{-g}\,F^{ik}\right)
  =-\frac{4\pi}{c}j^i
\end{equation}
（使用了公式（86.10)）.

最后，引力场和电磁场中的带电粒子的运动方程可通过将（23.4)中的四维加速度 $\mathrm{d}u^i/\mathrm{d}s$ 替换为 $\mathrm{D}u^i/\mathrm{d}s$ 来获得：
\begin{equation}
  mc\frac{\mathrm{D}u^i}{\mathrm{d}s}
  =mc\left(\frac{\mathrm{d}u^i}{\mathrm{d}s}
  +\Gamma^i_{kl}u^ku^l\right)
  =\frac{e}{c}F^{ik}u_k.
\end{equation}

\begin{xiti}

\noindent{\heiti 习题}\quad 按照如下定义引入三维矢量 $\boldsymbol{E}$，$\boldsymbol{D}$ 和反对称三维张量 $B_{\alpha\beta}$ 和 $H_{\alpha\beta}$：
\begin{equation*}
  E_{\alpha}=F_{0\alpha},\qquad B_{\alpha\beta}=F_{\alpha\beta},
\end{equation*}
\begin{equation*}
  D^{\alpha}=-\sqrt{g_{00}}\,F^{0\alpha},\qquad
  H^{\alpha\beta}=\sqrt{g_{00}}\,F^{\alpha\beta}.
\tag{1}\end{equation*}
写出给定引力场中的三维形式的麦克斯韦方程组（在带有度规 $\gamma_{\alpha\beta}$ 的三维空间中）.

解：上面引入的各量是不独立的.写出方程式
\begin{equation*}
  F_{0\alpha}=g_{0l}g_{\alpha m}F^{lm},\qquad
  F^{\alpha\beta}=g^{\alpha l}g^{\beta m}F_{lm},
\end{equation*}
并引入三维度规张量 $\gamma_{\alpha\beta}=-g_{\alpha\beta}+hg_{\alpha}g_{\beta}$（g 和 h 来自(88.11)），并应用公式（84.9）和（84.12），我们得到：
\begin{equation*}
  D_{\alpha}=\frac{E_{\alpha}}{\sqrt{h}}+g^{\beta}H_{\alpha\beta},
  \qquad
  B^{\alpha\beta}=\frac{H^{\alpha\beta}}{\sqrt{h}}
  +g^{\beta}E^{\alpha}-g^{\alpha}E^{\beta}.
\tag{2}\end{equation*}
引入矢量 $\boldsymbol{B}$，$\boldsymbol{H}$，它们分别与张量 $B_{\alpha\beta}$ 和 $H_{\alpha\beta}$ 对偶，按照定义
\begin{equation*}
  B^{\alpha}=-\frac{1}{2\sqrt{\gamma}}e^{\alpha\beta\gamma}B_{\beta\gamma},
  \qquad
  H_{\alpha}=-\frac{1}{2}\sqrt{\gamma}\,e_{\alpha\beta\gamma}H^{\beta\gamma}
\tag{3}\end{equation*}
（对比§88习题1的脚注)；引入负号的目的是为了使得伽利略坐标系中的矢量 H 和 B 等同于通常的磁场强度）.那么(2)式可以写成下面的形式
\begin{equation*}
  \boldsymbol{D}=\frac{\boldsymbol{E}}{\sqrt{h}}
  +\boldsymbol{H}\times\boldsymbol{g},\qquad
  \boldsymbol{B}=\frac{\boldsymbol{H}}{\sqrt{h}}
  +\boldsymbol{g}\times\boldsymbol{E}.
\tag{4}\end{equation*}
将定义 (1) 引入(90.2)式，我们得到方程
\begin{equation*}
  \frac{\partial B_{\alpha\beta}}{\partial x^{\gamma}}
  +\frac{\partial B_{\gamma\alpha}}{\partial x^{\beta}}
  +\frac{\partial B_{\beta\gamma}}{\partial x^{\alpha}}=0,
\end{equation*}
\begin{equation*}
  \frac{\partial B_{\alpha\beta}}{\partial x^0}
  +\frac{\partial E_{\alpha}}{\partial x^{\beta}}
  -\frac{\partial E_{\beta}}{\partial x^{\alpha}}=0,
\end{equation*}
或者，过渡到 (3) 式中的对偶量
\begin{equation*}
  \operatorname{div}\boldsymbol{B}=0,\qquad
  \operatorname{rot}\boldsymbol{E}
  =-\frac{1}{c\sqrt{\gamma}}\frac{\partial}{\partial t}
  \left(\sqrt{\gamma}\,\boldsymbol{B}\right)
\tag{5}\end{equation*}
（$x^0=ct$；算符 rot 和 div 的定义见§88习题 1的脚注).类似地，我们从(90.6)中得到方程
\begin{equation*}
  \frac{1}{\sqrt{\gamma}}\frac{\partial}{\partial x^{\alpha}}
  \left(\sqrt{\gamma}\,D^{\alpha}\right)=4\pi\rho,
\end{equation*}
\begin{equation*}
  \frac{1}{\sqrt{\gamma}}\frac{\partial}{\partial x^{\beta}}
  \left(\sqrt{\gamma}\,H^{\alpha\beta}\right)
  +\frac{1}{\sqrt{\gamma}}\frac{\partial}{\partial x^0}
  \left(\sqrt{\gamma}\,D^{\alpha}\right)
  =-4\pi\rho\frac{\mathrm{d}x^{\alpha}}{\mathrm{d}x^0},
\end{equation*}
或者用三维的矢量表示：
\begin{equation*}
  \operatorname{div}\boldsymbol{D}=4\pi\rho,\qquad
  \operatorname{rot}\boldsymbol{H}
  =\frac{1}{c\sqrt{\gamma}}\frac{\partial}{\partial t}
  \left(\sqrt{\gamma}\,\boldsymbol{D}\right)
  +\frac{4\pi}{c}\boldsymbol{s},
\tag{6}\end{equation*}
其中 $\boldsymbol{s}$ 是分量为 $s^{\alpha}=\rho\,\mathrm{d}x^{\alpha}/\mathrm{d}t$ 的矢量.

我们同样将连续性方程（90.5）写成三维形式
\begin{equation*}
  \frac{1}{\sqrt{\gamma}}\frac{\partial}{\partial t}
  \left(\sqrt{\gamma}\,\rho\right)
  +\operatorname{div}\boldsymbol{s}=0.
\tag{7}\end{equation*}
读者应注意到方程 (5)，(6) 与物质介质中电磁场的麦克斯韦方程组的相似性(当然，纯粹是形式上的).特别地，在静态引力场中，包含对时间的导数的各项中的 $\sqrt{\gamma}$ 会消失，而关系式(4)化简为 $\boldsymbol{D}=\boldsymbol{E}/\sqrt{h}$，$\boldsymbol{B}=\boldsymbol{H}/\sqrt{h}$.可以说，在对电磁场的影响这一方面，静态引力场起着介电常数和磁导率为 $\varepsilon=\mu=1/\sqrt{h}$ 的介质的作用.

\end{xiti}
